\section{Hodge numbers, cotangent direct images, and the Fr\"olicher page}
\label{sec:hodge-frolicher}

This section first derives the one-parameter Hodge table and the complete
Fr\"olicher spectral sequence from the candidate-specific structure-sheaf
calculation and the verified sphere topology.  It then performs the missing
logarithmic direct-image calculation through all three special fibres.  The
cyclic invariant lattices, logarithmic Kodaira--Spencer divisor, residue
cancellation, and one remaining global scalar give
$h^{1,2}=1$ and $h^{1,1}=2$; thus no Hodge number remains undetermined.

\subsection{Vanishing of holomorphic forms}

We first place the holomorphic-form calculation inside the workbench rather
than using it as an unexpanded assertion from the reconstructed source.

\begin{proposition}[candidate-specific holomorphic-form vanishing]
\label{prop:workbench-holomorphic-form-vanishing}
\Verified\ For the constructed threefold,
\[
 H^0(X,\Omega_X^p)=0\qquad (p=1,2,3).
\]
\end{proposition}

\begin{proof}
Over the marked cover of the smooth part, write the fibre coordinates as
$\zeta=(\zeta_1,\zeta_2)$ and
\[
 \Pi=[Z\mid I_2],\qquad
 Z=\begin{pmatrix}6\mu&\tau\\ \beta&\mu\end{pmatrix}.
\]
Translation by $\lambda\in\Lambda$ sends
$d\zeta$ to $d\zeta+\dot\Pi\lambda\,dz$.  Compactness of each torus fibre
therefore puts every translation-invariant holomorphic form into one of the
following normal forms:
\begin{align*}
 p=1:&\quad a(z)\,dz+c(z)\cdot d\zeta,
       &c(z)\dot Z(z)&=0,\\
 p=2:&\quad dz\wedge\bigl(b(z)\cdot d\zeta\bigr),\\
 p=3:&\quad c(z)\,dz\wedge d\zeta_1\wedge d\zeta_2.
\end{align*}
Indeed, fibrewise periodic coefficients are constant.  A possible affine
part in a coefficient vanishes after restriction to the real span of the
two constant columns of $\Pi$; in degree two this also kills the putative
$d\zeta_1\wedge d\zeta_2$ term because $\dot\tau$ is nonzero on the smooth
base.

The finite fillings are obtained from smooth torus families by free cyclic
quotients, so the displayed coefficients extend holomorphically across the
two elliptic lifts.  At the cusp use the verified toric charts
$t=w_0w_1w_2$.  Substitution of
\[
 d\zeta_i=(2\pi i)^{-1}dx_i/x_i,
 \qquad ds=(2\pi i)^{-1}dt/t
\]
and restriction successively to the coordinate lines and planes shows that
the coefficients of $d\zeta_i$, $ds$, $ds\wedge d\zeta_i$, and
$ds\wedge d\zeta_1\wedge d\zeta_2$ vanish to order at least one in $t$.
Indeed, on every unimodular monomial chart the triple
$(d\log x_1,d\log x_2,d\log t)$ is a constant invertible linear combination
of the $d\log w_j$.  A coefficient $F(t)$ multiplying a nonzero logarithmic
residue is holomorphic only if its constant term vanishes: restricting to
$w_j=0$ gives that residue as $F(0)$, so $F(t)=tG(t)$.  Invertibility of the
monomial change of basis applies this argument to each displayed
coefficient.
Thus every coefficient used below descends with the asserted regularity at
all three points of the base.

The linear part of a deck transformation has the form
\[
 R_g=\begin{pmatrix}r_g&0\\ *&1\end{pmatrix},
 \qquad \jmath_g:=r_g^{-1}=c_g\tau+d_g,
 \qquad g^*(d\tau)=\jmath_g^{-2}d\tau.
\]
For a one-form, the second component $c_2$ is a holomorphic invariant
function on the base and vanishes at the cusp, hence is zero.  The second
column of $c\dot Z=0$ then gives $c_1\dot\tau=0$, so $c_1=0$.  What remains
is a holomorphic one-form on $\mathbb P^1$, and hence vanishes.

For a two-form, $b_2\,dz$ descends to a holomorphic one-form on
$\mathbb P^1$ and is zero.  The remaining form
$\psi=b_1\,dz$ satisfies $g^*\psi=\jmath_g\psi$.  Consequently
\[
 \psi^{\otimes2}\otimes d\tau
\]
is invariant.  At the order-three and order-four points it descends as a
cubic differential with a pole of order at most two, while the toric cusp
estimate gives a pole of order at most one.  It is therefore a section of
\[
 (\Omega^1_{\mathbb P^1})^{\otimes3}
 (2p_1+2p_2+p_0)\simeq\mathcal O_{\mathbb P^1}(-1),
\]
so it is zero.  Since $\dot\tau$ is nonzero on a dense open set,
$\psi=0$.  The coefficient of a three-form obeys the same weight-one law
and the same argument applies.  Hence all three spaces vanish.
\end{proof}

\subsection{The exact one-parameter Hodge diamond}

\begin{theorem}[Hodge diamond determined up to one integer]
\label{thm:parametric-hodge-diamond}
\Verified\ Put $m=h^{1,2}(X)$.  Then $m\in\mathbb Z_{\ge0}$ and the
only nonzero Hodge numbers are
\[
 h^{0,0}=h^{0,1}=h^{3,2}=h^{3,3}=1,
 \qquad
 h^{1,1}=h^{2,2}=m+1,
 \qquad
 h^{1,2}=h^{2,1}=m.
\]
Equivalently,
\[
\begin{array}{c|cccc}
q\backslash p&0&1&2&3\\ \hline
3&0&0&0&1\\
2&0&m&m+1&1\\
1&1&m+1&m&0\\
0&1&0&0&0
\end{array}
\]
where the row index is $q$ and the column index is $p$.
\end{theorem}

\begin{proof}
The independently reconstructed structure-sheaf direct images are
\[
 f_*\mathcal O_X=\mathcal O,
 \qquad R^1f_*\mathcal O_X=\mathcal O\oplus\mathcal O(-1),
 \qquad R^2f_*\mathcal O_X=\mathcal O(-1)
\]
by \cref{thm:structure-direct-images-audit}.  Coherent Leray on
$\mathbb P^1$ therefore gives
\[
 (h^{0,0},h^{0,1},h^{0,2},h^{0,3})=(1,1,0,0).
\]
Together with \cref{prop:workbench-holomorphic-form-vanishing} and Serre
duality \cite[\S13, Theorem~4, p.~22]{Serre1955Duality}, this fixes every
entry except the four middle entries, with
$h^{1,2}=h^{2,1}=m$ and $h^{1,1}=h^{2,2}$.

The verified integral sphere topology gives $c_1(X)=c_2(X)=0$ and
$\int_Xc_3(X)=2$.  In complex dimension three,
\begin{align*}
 \operatorname{td}(X)
 &=1+\frac12c_1+\frac1{12}(c_1^2+c_2)+\frac1{24}c_1c_2,\\
 \operatorname{ch}(\Omega_X^1)
 &=3-c_1+\frac12(c_1^2-2c_2)
   -\frac16(c_1^3-3c_1c_2+3c_3).
\end{align*}
Hirzebruch--Riemann--Roch for arbitrary compact complex manifolds
\cite[Theorem~3, p.~427]{AtiyahSinger1963} hence gives
\[
 \chi(\Omega_X^1)=-\frac12\int_Xc_3=-1.
\]
The already known zero entries reduce this identity to
$-h^{1,1}+h^{1,2}=-1$, proving $h^{1,1}=m+1$.
\end{proof}

This specializes the general constraints of Gray and Ugarte; a convenient
published synthesis is Angella \cite{Angella2018}, while the second-page
symmetry and inequalities originate in \cite{Ugarte2000}.  The argument
above uses the stronger candidate-specific direct images and form
vanishings, rather than assuming a general Hodge pattern for a hypothetical
complex $S^6$.

\begin{theorem}[the complete Fr\"olicher page]
\label{thm:frolicher-e2-candidate}
\Verified\ The Fr\"olicher spectral sequence of $X$ does not degenerate at
$E_1$, but it does degenerate at $E_2$.  Apart from the endpoint copies of
$\mathbb C$, the two nonzero $d_1=\partial$ rows are exact complexes
\[
 0\longrightarrow\mathbb C
 \longrightarrow\mathbb C^{m+1}
 \longrightarrow\mathbb C^m
 \longrightarrow0,
\]
\[
 0\longrightarrow\mathbb C^m
 \longrightarrow\mathbb C^{m+1}
 \longrightarrow\mathbb C
 \longrightarrow0.
\]
Their differential ranks are respectively $(1,m)$ and $(m,1)$.
Consequently
\[
 \dim_{\mathbb C}E_1=6+4m,
 \qquad \dim_{\mathbb C}E_2=2,
 \qquad E_2=E_\infty.
\]
\end{theorem}

\begin{proof}
Fr\"olicher's original construction identifies $E_1$ with Dolbeault
cohomology and $E_\infty$ with the associated graded de Rham cohomology
\cite[p.~643, Lemmas~1--2 and Theorem~3]{Frolicher1955}.  The table in
\cref{thm:parametric-hodge-diamond} leaves only the displayed
$q=1$ and $q=2$ row complexes between the endpoint classes.  Since
$H^k_{\mathrm{dR}}(X;\mathbb C)=0$ for $1\le k\le5$, every middle
$E_\infty^{p,q}$ is zero.  There is only one a priori possible higher
differential between middle positions,
$d_2:E_2^{0,2}\to E_2^{2,1}$, but already
$E_1^{0,2}=H^2(X,\mathcal O_X)=0$.  Every other $d_r$ with $r\ge2$ has
zero source or target by the same Hodge table.  Thus $E_2=E_\infty$, so the
two rows must be exact.  The ranks follow from their dimensions.
\end{proof}

The theorem explains why neither topology nor the Fr\"olicher sequence alone
can select a value of $m$: the two exact row complexes exist numerically for
every $m\ge0$.  The singular-fibre calculation below supplies the missing
analytic information.

\subsection{Generic maximal variation and its elliptic zero}

The most tempting shortcut to the remaining Hodge number is to treat the
Kodaira--Spencer matrix as invertible everywhere.  The following calculation
shows precisely what is true and what is false.

\begin{proposition}[determinant of the printed period derivative]
\label{prop:period-derivative-determinant}
\Verified\ With dots denoting differentiation in the coordinate on
$\mathfrak h_z$, set
\[
 K=\det(\dot Z)=6\dot\mu^{\,2}-\dot\tau\dot\beta.
\]
Then $K$ is not identically zero, so the Kodaira--Spencer map has rank two
on a dense open set.  However $K(z_1)=0$ at the chosen order-three elliptic
point.  In particular, the assertion $\dot Z\ne0$ does not imply that
$\dot Z$ is invertible, and direct images across the finite fibres require
their actual cyclic invariant lattices.
\end{proposition}

\begin{proof}
At the cusp, \cref{prop:period-input-audit} gives
\[
 t=e^{2\pi i s},\qquad Z=sB_0+C(t),\qquad
 B_0=\begin{pmatrix}0&1\\-1&0\end{pmatrix},
\]
with $C$ holomorphic at $t=0$.  Hence
\[
 \frac{dZ}{ds}=B_0+2\pi i\,tC'(t),
 \qquad
 \det\!\left(\frac{dZ}{ds}\right)=1+O(t).
\]
The cusp coordinate $s$ is locally biholomorphic on the marked cover, so
$K$ is nonzero near the punctured cusp and therefore is not identically
zero.

Now put $\rho=e^{\pi i/3}=\tau(z_1)$.  The difference between $\mu$ and
the local affine solution
\[
 \mu_1=\frac{2-\tau}{3}
\]
vanishes to order at least two at $z_1$: the homogeneous line has local
generator of order two there.  Since $\tau-\rho$ has a simple zero,
\[
 \dot\mu(z_1)=-\frac13\dot\tau(z_1).
\]
Write
\[
 \phi_1=2-\frac{6(1-\mu)^2}{\tau},
 \qquad \beta\circ g_1=\beta+\phi_1.
\]
Using $\rho^2-\rho+1=0$ and the preceding derivative gives
\[
 \dot\phi_1(z_1)
 =-\frac23(1+\rho)\dot\tau(z_1).
\]
The derivative of the linearized order-three action at its fixed point is
$g_1'(z_1)=-\rho$.  Differentiating the transformation law for $\beta$
therefore yields
\[
 -(1+\rho)\dot\beta(z_1)=\dot\phi_1(z_1),
 \qquad
 \dot\beta(z_1)=\frac23\dot\tau(z_1).
\]
Substitution gives
\[
 K(z_1)=6\left(-\frac13\dot\tau(z_1)\right)^2
 -\dot\tau(z_1)\left(\frac23\dot\tau(z_1)\right)=0.
\]
\end{proof}

The zero divisor of $K$ cannot be read from the single equality $K(z_1)=0$.
It is determined below only after the invariant lattices at the two multiple
fibres have been inserted.  This distinction is essential: direct image does
not commute with taking the dual lattice at a ramified coarse fibre.

\subsection{Logarithmic cotangent direct images}

Put
\[
 B=\mathbb P^1,\qquad
 \Sigma=p_0+p_1+p_2,\qquad
 D=f^{-1}(\Sigma)_{\mathrm{red}}=W+S_1+S_2,
\]
and set
\[
 \mathcal L=\Omega_B^1(\log\Sigma)\simeq\mathcal O_B(1),\qquad
 \mathcal A=\Omega_X^1(\log D),\qquad
 \mathcal E=\mathcal A/f^*\mathcal L.
\tag{11.1}
\]
At the cusp the local equations $t=z_1\cdots z_r$, $1\le r\le3$, show
directly that $\mathcal A$ and $\mathcal E$ are locally free.  At a finite
fibre the covering family has divisor $s=0$, the cyclic action is free on the
total space, and $dt/t=m\,ds/s$; hence its locally free logarithmic sequence
descends through the finite \`etale quotient.  Thus everywhere one has the
logarithmic conormal sequence
\[
 0\longrightarrow f^*\mathcal L\longrightarrow\mathcal A
 \longrightarrow\mathcal E\longrightarrow0.
\tag{11.2}
\]
The local $d\log$ bases and their relative form, including
$t=\prod z_i^{e_i}$ and analytic base change, are the logarithmic conormal
formalism of \cite[Chapter~II, \S3.2--3.3, pp.~72--75]{Deligne1970}; the
display above also follows directly from the stated local equations.

\begin{proposition}[the three logarithmic relative direct images]
\label{prop:log-relative-direct-images}
\Verified\ One has
\begin{align*}
 f_*\mathcal E&\simeq\mathcal O\oplus\mathcal O(-1),\\
 R^1f_*\mathcal E&\simeq
       \mathcal O^{\oplus2}\oplus\mathcal O(-1)^{\oplus2},\\
 R^2f_*\mathcal E&\simeq\mathcal O\oplus\mathcal O(-1).
\end{align*}
All three sheaves are locally free and commute with base change.
\end{proposition}

\begin{proof}
At the cusp, the toric coordinates give
\[
 \mathcal E|_{N_0}
 =\mathcal O_{N_0}[d\log x_1]\oplus
  \mathcal O_{N_0}[d\log x_2].
\tag{11.3}
\]
Indeed a quotient shear changes each $d\log x_i$ only by a holomorphic
multiple of $d\log t$, which is zero in $\mathcal E$.  The already proved
structure-sheaf calculation therefore gives free local direct images of
ranks $2,4,2$ at $p_0$.

We now keep the ramification at $p_j$.  Write
\[
 A=\mathbb C\{t\},\qquad C=\mathbb C\{s\},\qquad t=s^m,
 \qquad (m,a)=(3,2)\text{ or }(4,1).
\]
Cartan's analytic quotient theorem identifies local holomorphic functions on
the quotient with invariants \cite[\S1, Theorem~1, pp.~1--2]{Cartan1954Quotient}.
The action on the total space is free, so the quotient map is finite \'etale;
the corresponding locally free sheaves and their equivariant structures
descend by finite \'etale descent.  This sheaf statement is not being
attributed to Cartan's theorem on invariant functions.  Finite pushforward
is exact, and Reynolds averaging
$v\mapsto |G_j|^{-1}\sum_{g\in G_j}gv$ makes the invariant functor exact in
characteristic zero.  Hence
\[
 R^qf_*\mathcal E|_{D_j}
 =\bigl(\pi_*R^q\widetilde f_*\widetilde{\mathcal E}\bigr)^{G_j}.
\tag{11.3a}
\]
Choose equivariant frames $e_0,e_1$ of the relative cotangent and
$b_0,b_1$ of $R^1\widetilde f_*\mathcal O$ whose semilinear characters are
respectively
\[
 (1,\zeta^a),\qquad (1,\zeta^{-a});
\]
$b_0$ is the nowhere-zero class $\gamma^{0,1}$.  These may be chosen as
holomorphic constant-character frames, rather than merely as eigenvectors on
the central fibre.  Indeed, if $w$ lifts a $\nu$-eigenvector at $s=0$, then
\[
 P_\nu(w)=\frac1m\sum_{k=0}^{m-1}\nu^{-k}(g^*)^kw
\]
satisfies $g^*P_\nu(w)=\nu P_\nu(w)$ and has the same nonzero central value;
after shrinking the disc, the projected eigenvectors form a frame.  On the
marked torus family the projection formula gives
\[
 R^q\widetilde f_*\widetilde{\mathcal E}
 \simeq
 (\mathcal O e_0\oplus\mathcal O e_1)
 \otimes
 \bigwedge^q(\mathcal O b_0\oplus\mathcal O b_1),
 \qquad 0\le q\le2.
\]
Solving
$\nu\zeta^r=1$ for the invariant exponent gives the following \emph{actual}
$A$-bases:
\begin{align*}
 f_*\mathcal E:\quad& e_0,\ s^{m-a}e_1,\\
 R^1f_*\mathcal E:\quad&
 e_0b_0,\ s^ae_0b_1,\ s^{m-a}e_1b_0,\ e_1b_1,\\
 R^2f_*\mathcal E:\quad&
 s^ae_0(b_0\wedge b_1),\ e_1(b_0\wedge b_1).
\tag{11.4}
\end{align*}
For comparison, the variable tangent lattice is generated by $s^av_1$.
Thus evaluation of the descended variable cotangent on the descended
variable tangent is $s^m=t$, not a unit.  This proves explicitly why
\[
 f_*\mathcal E\longrightarrow(f_*T_{X/B})^\vee
\]
is multiplication by $\operatorname{diag}(1,t)$ locally, although
$\mathcal E^\vee\simeq T_{X/B}$ before direct image.  This last identity is
an identity of locally free relative sheaves, not an assertion that formation
of the kernel commutes with taking a tangent-space fibre.  For example, at a
triple chart
\[
 R=\mathbb C\{x,y,z\},\qquad t=xyz,
\]
one has
\[
 T_{X/B}=\ker\!\left[R^3\xrightarrow{(yz,xz,xy)}R\right].
\]
If $yz\,u+xz\,v+xy\,w=0$, reduction modulo $x$ and $y$ gives
$u=xU$ and $v=yV$, and cancellation in the domain gives
$w=-z(U+V)$.  Hence
\[
 T_{X/B}
 =R(x\partial_x-z\partial_z)\oplus
  R(y\partial_y-z\partial_z)
 =\mathcal E^\vee.
\]
The kernel of the pointwise linear map $df_x$ at the triple point is
three-dimensional, whereas the fibre of this sheaf kernel is
two-dimensional: tensoring the kernel sequence with the residue field is
not left exact because the image ideal $(xy,xz,yz)$ is not flat.  Thus a
section can be nowhere zero in $T_{X/B}=\mathcal E^\vee$ even when its image
in the ambient tangent-space fibre vanishes.

The invariant vertical field $\xi$ gives a trivial subline of
$V:=f_*T_{X/B}$; the quotient has the reciprocal automorphy of the Hodge
line $\mathfrak K\simeq\mathcal O(1)$ from
\cref{lem:hodge-line-audit}.  Since $H^1(\mathcal O(1))=0$, the resulting
extension splits:
\[
 V\simeq\mathcal O\xi\oplus\mathcal O(-1),
 \qquad V^\vee\simeq\mathcal O\oplus\mathcal O(1).
\]
Evaluation gives a generically invertible map
$f_*\mathcal E\to V^\vee$.  Formula (11.4) shows that its elementary
divisors are $(1,t)$ at each of $p_1,p_2$ and $(1,1)$ elsewhere.  Thus its
cokernel has total length two and
\(\deg(f_*\mathcal E)=1-2=-1\).  Evaluation on $\xi$ is stalkwise
surjective: the covector dual to the fixed eigenvector is a unit generator
at the cusp and in both finite invariant bases.  Its kernel is consequently
a line bundle of degree $-1$, and
\[
 0\longrightarrow\mathcal O(-1)\longrightarrow f_*\mathcal E
 \xrightarrow{\operatorname{ev}_\xi}\mathcal O\longrightarrow0.
\]
Because $H^1(\mathcal O(-1))=0$, this extension also splits, proving
\[
 f_*\mathcal E\simeq\mathcal O\oplus\mathcal O(-1).
\]
Also, by 
\cref{thm:structure-direct-images-audit},
\[
 G:=R^1f_*\mathcal O_X=\mathcal O\oplus\mathcal O(-1),
 \qquad R^2f_*\mathcal O_X=\det G=\mathcal O(-1).
\tag{11.5}
\]
Here is the global filtration implicit in those bases.  Choose the
degree-zero summand in the preceding split sequence so that its generator
evaluates to one on $\xi$; locally this is the fixed eigen-covector $e_0$.
Together with the fixed class $b_0=\gamma^{0,1}$, take the corresponding
saturated sublines in $f_*\mathcal E$ and $G$.  Formula (11.3) at the cusp
and the bases (11.4) at the finite points show that these sublines are locally
generated by $e_0,b_0$ and that their quotients are locally generated by
$s^{m-a}e_1,s^ab_1$; hence both inclusions and both quotients are locally
free.  The evaluation splitting and Hodge-line calculation identify the two
filtrations as
\[
 0\subset\mathcal O\subset f_*\mathcal E,\qquad
 0\subset\mathcal O\subset G,\qquad
 \operatorname{gr}(f_*\mathcal E)=
 \operatorname{gr}(G)=\mathcal O\oplus\mathcal O(-1).
\]
Over the smooth locus the cup-product identification is
$R^1f_*\mathcal E=(f_*\mathcal E)\otimes G$.  Formula (11.4) extends its four
character lines as saturated rank-one subsheaves over all of $B$.  Three
associated-graded lines are generated locally by
$e_0b_0,s^ae_0b_1,s^{m-a}e_1b_0$ and have degrees $0,-1,-1$.  On the fourth,
the separately descended factors would
give
\[
 (s^{m-a}e_1)(s^ab_1)=t(e_1b_1),
\]
whereas $e_1b_1$ itself is invariant.  Thus that graded line is the upper
modification of $\mathcal O(-2)$ at $p_1+p_2$, hence is $\mathcal O$.
Similarly, the two graded lines of $R^2f_*\mathcal E$ are
$\mathcal O(-1)$ and
$\mathcal O(-2)(p_1+p_2)=\mathcal O$.  Every extension between the resulting
degree-zero and degree-minus-one lines splits: its $\operatorname{Ext}^1$
is one of $H^1(\mathcal O(-1))$, $H^1(\mathcal O)$, or
$H^1(\mathcal O(1))$, all zero.  This proves the three asserted splittings.
This is also the elementary filtered form of the
Birkhoff--Grothendieck classification
\cite[Theorem~2.1, p.~126]{Grothendieck1957}; the divisor modifications and
$\operatorname{Ext}^1$ identifications may equivalently be read in
\cite[Tags 0C6K, 0B3A, and 01XT]{StacksProject}.

Finally, (11.3a)--(11.4) make every direct image free of ranks $2,4,2$ at
the finite points; (11.3) does the same at the cusp, and the smooth locus is
standard.  The multiple-fibre base change can also be read directly from the
invariant modules.  The Reynolds idempotent is $A$-linear, so invariants
commute with arbitrary finite $A$-base change.  For the scheme-theoretic
special fibre,
\[
 C\otimes_AA/(t)=C/(s^m),
 \qquad
 \bigl((C/(s^m))v\bigr)^G=\mathbb C\,s^rv
\]
on a character line whose invariant exponent satisfies $0\le r<m$.
Every basis vector in (11.4) therefore survives reduction modulo $t$, and
the special-fibre dimensions are exactly $2,4,2$.  No $R^q$ with $q>2$
occurs.  Descending induction in analytic cohomology and base change then
makes every base-change map an isomorphism, including for the nonreduced
finite fibres; precisely the same criterion, with citation, is recorded in
the proof of
\cref{thm:structure-direct-images-audit}; see in particular
\cite[Chapter~III, \S3, Theorem~3.4 and Corollary~3.10]{BanicaStanasila1976}.
\end{proof}

\begin{proposition}[the logarithmic Kodaira--Spencer divisor]
\label{prop:log-ks-divisor}
\Verified\ The connecting maps of (11.2) are
\begin{align*}
 \kappa_0&:f_*\mathcal E\longrightarrow
  \mathcal L\otimes R^1f_*\mathcal O_X,\\
 \kappa_1&:R^1f_*\mathcal E\longrightarrow
  \mathcal L\otimes R^2f_*\mathcal O_X\simeq\mathcal O.
\end{align*}
Then
\[
 \operatorname{div}(\det\kappa_0)=p_1+p_2,
 \qquad \kappa_1\text{ is surjective as a sheaf}.
\tag{11.6}
\]
In particular $\kappa_0$ is injective and
$\operatorname{coker}\kappa_0=\mathbb C_{p_1}\oplus\mathbb C_{p_2}$.
Moreover the raw period determinant has
\[
 \operatorname{ord}_{z_1}K=2,\qquad
 \operatorname{ord}_{z_2}K=0.
\tag{11.7}
\]
\end{proposition}

\begin{proof}
On the smooth locus, the connecting construction and its
exterior-product Leibniz rule are the Gauss--Manin/Kodaira--Spencer
formalism of
\cite[\S2, equations~(5)--(10), pp.~202--204]{KatzOda1968}.
Katz--Oda work with a smooth algebraic morphism; the citation is not being
used as a logarithmic analytic special-fibre theorem.  In the present family,
Deligne's local monomial calculation supplies the locally split logarithmic
sequence (11.2) \cite[Chapter~II, \S3.3]{Deligne1970}.  Taking a local
splitting and a \v Cech representative, the connecting cochain is the
difference of the two splittings.  Since exterior product of cochains obeys
\[
 \delta(\alpha\wedge\beta)
 =\delta\alpha\wedge\beta+(-1)^{\deg\alpha}
   \alpha\wedge\delta\beta,
\]
the two displayed connecting maps extend across every logarithmic chart and
satisfy the same Leibniz rule.  Thus the singular-fibre formula below is a
direct calculation from (11.2), not an extrapolation of the Katz--Oda
hypotheses.
The determinant is a section of
\[
 \det(\mathcal L\otimes G)\otimes\det(f_*\mathcal E)^{-1}
 =\mathcal L^{\otimes2}\otimes\det G
   \otimes\det(f_*\mathcal E)^{-1}
 \simeq\mathcal O(2)\otimes\mathcal O(-1)\otimes\mathcal O(1)
 \simeq\mathcal O(2).
\tag{11.8}
\]
At the cusp, $Z=sB_0+C(t)$ and $dZ/ds=B_0+O(t)$, so it is a unit.  At a
finite point use $\ell=dt/t=m\,ds/s$.  In invariant determinant frames
with no normalization imposed against the raw marked frames, direct
calculation gives
\[
 \det\kappa_0
 =\varepsilon(t)m^{-2}s^{m+2-2a}K(s),
 \qquad \varepsilon(0)\ne0.
\tag{11.9}
\]
Here $\varepsilon$ is the nowhere-zero invariant holomorphic unit recording the two
frame normalizations; compatible eigenframes may be normalized so that
$\varepsilon=1$.  The dot in $K$ means $d/ds$ in the chosen linearizing coordinate.
Replacing $s$ by another such coordinate changes only the unit, so every
divisor and order calculation below is coordinate-independent.
Equivariance gives
\[
 \operatorname{ord}_sK\equiv2a-2\pmod m.
\tag{11.10}
\]
At $(m,a)=(3,2)$, 
\cref{prop:period-derivative-determinant} supplies $K(z_1)=0$; hence the
right side of (11.9) has coarse order at least one.  At $(m,a)=(4,1)$,
the prefactor is $s^4=t$, so the coarse order is again at least one.
The two forced zeros exhaust the degree in (11.8).  They are therefore
simple, there are no others, and (11.7) follows from (11.9)--(11.10).

For completeness, we derive the coarse matrices rather than only recording
their shape.  Let
\[
 D_{\mathrm{KS}}(s)=(d_{rs}(s))_{0\leq r,s\leq1}
\]
be the Kodaira--Spencer matrix on the marked cover after the common factor
$ds$ has been removed.  Since $t=s^m$ and
$ds=(s/m)\ell$, passage from the marked frames to the actual invariant
frames in (11.4) gives
\[
 A=\frac{s}{m}d_{00},\qquad
 B=\frac{s^{1-a}}m d_{01},\qquad
 C=\frac{s^{1-a}}m d_{10},\qquad
 D_0=\frac{s^{m+1-2a}}m d_{11}.
\tag{11.10a}
\]
The character of each expression in (11.10a) is trivial, so each is a
holomorphic function of $t$.  In the bases
$(e_0,s^{m-a}e_1)$ and $(b_0,s^ab_1)$ one consequently has the exact matrix
\[
 [\kappa_0]=\begin{pmatrix}A&tB\\ C&D_0\end{pmatrix}.
\tag{11.10b}
\]
Its determinant is (11.9).  Applying the exterior-product Leibniz rule to
the four ordered generators
\[
 e_0b_0,\quad s^ae_0b_1,\quad
 s^{m-a}e_1b_0,\quad e_1b_1
\]
now gives, without a change of basis,
\[
 [\kappa_1]=(-C,\ A,\ -D_0,\ B)
\tag{11.11}
\]
into the invariant generator $s^a(b_0\wedge b_1)\ell$.  If (11.11) were
not surjective at a finite point, all four entries would lie in $(t)$ and
$\det\kappa_0$ would have order at least two, contrary to (11.6).  Away
from $p_1,p_2$, the determinant divisor (11.6) makes $\kappa_0$
invertible, so the exterior-product row (11.11) contains a unit and is
surjective there as well.  At the cusp this reduction is not an additional
choice.  Use the logarithmic base
frame
\(\ell=ds=(2\pi i)^{-1}dt/t\), the relative cotangent frame
\((e_0,e_1)=([d\log x_1],[d\log x_2])\) of (11.3), and the quotient frame
\((\gamma,u)\) of \(G\), which extends across the cusp by
\cref{thm:structure-direct-images-audit}.  In these frames the smooth-locus Kodaira--Spencer matrix is the
transpose of the derivative of the period matrix: for
\(h=a e_0+b e_1\),
\[
 \kappa_0(h)=
 (6a\dot\mu+b\dot\beta)\gamma+
 (a\dot\tau+b\dot\mu)u .
 \tag{11.11a}
\]
Here the dot is \(d/ds\).  The cusp expansion
\(Z=sB_0+C(t)\), \(t=e^{2\pi i s}\), gives
\[
 \frac{dZ}{ds}=B_0+2\pi i\,tC'(t),
 \qquad
 B_0=\begin{pmatrix}0&1\\-1&0\end{pmatrix}.
\]
These logarithmic frames extend over \(t=0\), so (11.11a) reduces exactly
to \(B_0^{\mathsf t}\).  Write the cusp matrix, separately from the
finite-fibre normalization above, as
\([\kappa_0]_{p_0}=\left(\begin{smallmatrix}A&E\\C&D\end{smallmatrix}\right)\).
In the same ordered exterior basis used for (11.11), its Leibniz row is
\((-C,A,-D,E)\).  Substituting
\((A,E,C,D)=(0,-1,1,0)\) gives the second reduction.  Thus, in the named
frames,
\[
 [\kappa_0]\equiv\begin{pmatrix}0&-1\\1&0\end{pmatrix}\pmod t,
 \qquad [\kappa_1]\equiv(-1,0,0,-1)\pmod t,
\]
exactly.  Under arbitrary holomorphic changes of source and target frames
the matrices change by invertible left and right factors; in particular
invertibility of \(\kappa_0\) and the existence of a unit entry in the row
of \(\kappa_1\) are unchanged.  Thus $\kappa_1$ is surjective everywhere.
\end{proof}

\subsection{Residues and the absolute cotangent direct images}

Let $i_D:D\hookrightarrow X$, and let
\[
 \nu:\widetilde D=\mathrm{dP}_6\sqcup S_1\sqcup S_2\longrightarrow D
\]
be the normalization.  In formulas below, $\nu_*\mathcal O_{\widetilde D}$
abbreviates $(i_D\circ\nu)_*\mathcal O_{\widetilde D}$ as a sheaf on $X$.
The ordinary/logarithmic local calculation for
$z_1\cdots z_r=0$ glues to
\[
 0\longrightarrow\Omega_X^1\longrightarrow\mathcal A
 \xrightarrow{\operatorname{res}}\nu_*\mathcal O_{\widetilde D}
 \longrightarrow0.
\tag{11.12}
\]
The normalization in the target is indispensable at the self-crossing
cusp.  This is the degree-one residue sequence associated with the normalized
strata of a normal-crossings divisor
\cite[Chapter~II, \S3.4--3.7, pp.~76--78]{Deligne1970}; the displayed local
calculation fixes its convention in the present self-crossing model.

\begin{proposition}[absolute cotangent direct images]
\label{prop:absolute-cotangent-direct-images}
\Verified\ Put $\mathcal K=\ker\kappa_1$.  Then
\[
 f_*\Omega_X^1\simeq\mathcal O(-2),\qquad
 R^1f_*\Omega_X^1\simeq\mathcal K,\qquad
 R^2f_*\Omega_X^1\simeq\mathcal O\oplus\mathcal O(-1).
\tag{11.13}
\]
None of these sheaves has torsion.
\end{proposition}

\begin{proof}
The three normalized components have
\[
 R^0f_*\nu_*\mathcal O_{\widetilde D}
  =\mathbb C_{p_0}\oplus\mathbb C_{p_1}\oplus\mathbb C_{p_2},
 \quad
 R^1f_*\nu_*\mathcal O_{\widetilde D}
  =\mathbb C_{p_1}\oplus\mathbb C_{p_2},
 \quad R^2=0.
\tag{11.14}
\]
Indeed $\mathrm{dP}_6$ is rational and $S_1,S_2$ are bielliptic, with
$h^1(\mathcal O_{S_j})=1$ and $h^2(\mathcal O_{S_j})=0$.

The long direct-image sequence of (11.2), together with
\cref{prop:log-ks-divisor}, gives
\[
 f_*\mathcal A=\mathcal O(1),\qquad
 0\longrightarrow T\longrightarrow R^1f_*\mathcal A
 \longrightarrow\mathcal K\longrightarrow0,
 \qquad R^2f_*\mathcal A=\mathcal O\oplus\mathcal O(-1),
\tag{11.15}
\]
where $T=\mathbb C_{p_1}\oplus\mathbb C_{p_2}$.
The residue of the base logarithmic form is weighted evaluation:
\[
 \mathcal O(1)\longrightarrow
 \mathbb C_{p_0}\oplus\mathbb C_{p_1}\oplus\mathbb C_{p_2},
 \qquad q\longmapsto(q(p_0),\,3q(p_1),\,4q(p_2)).
\tag{11.16}
\]
It is surjective as a sheaf and has kernel
$\mathcal O(1)(-p_0-p_1-p_2)=\mathcal O(-2)$.

It remains to identify the torsion in (11.15).  At $p_j$, the residue map on
$G\otimes\mathbb C_{p_j}$ sends
\[
 x b_0+y s^ab_1\longmapsto m_jx[b_0]
 \in H^1(S_j,\mathcal O_{S_j}).
\tag{11.17}
\]
More precisely, after trivializing $\mathcal L$ by
$\ell=dt/t$, the logarithmic conormal and residue maps form the exact
commutative diagram
\[
\begin{array}{ccccccccc}
0&\longrightarrow&f^*\mathcal L&\longrightarrow&\mathcal A
 &\longrightarrow&\mathcal E&\longrightarrow&0\\
&&\big\downarrow\,{\scriptstyle\operatorname{res}_B}
&&\big\downarrow\,{\scriptstyle\operatorname{res}_{S_j}}
&&\big\downarrow\,{\scriptstyle 0}&&\\
0&\longrightarrow&i_{j*}\mathcal O_{S_j}&\xrightarrow{\ =\ }
 i_{j*}\mathcal O_{S_j}&\longrightarrow&0&&
\end{array}
\tag{11.16a}
\]
with
\[
 \operatorname{res}_B(q\ell)=m_jq|_{S_j}.
\tag{11.16b}
\]
Applying $Rf_*$ to (11.16a) gives the typed composite
\[
 f_*\mathcal E|_{p_j}\xrightarrow{\ \kappa_0\ }
 (\mathcal L\otimes G)|_{p_j}
 \xrightarrow{\ \operatorname{res}_{1,j}\ }
 H^1(S_j,\mathcal O_{S_j}),
 \qquad \operatorname{res}_{1,j}\kappa_0=0,
\tag{11.16c}
\]
and (11.16b) gives exactly
\[
 \operatorname{res}_{1,j}(xb_0+ys^ab_1)=m_jx[b_0].
\]
Thus its kernel is the one-dimensional line
$\mathbb C\,s^ab_1$.  The simple zero of $\det\kappa_0$ proved in
\cref{prop:log-ks-divisor} implies that $\kappa_0\bmod t$ has rank one;
the zero composite in (11.16c) therefore forces
\[
 \operatorname{im}(\kappa_0\bmod t)
 =\ker\operatorname{res}_{1,j}
 =\mathbb C\,s^ab_1.
\tag{$11.17a_{\mathrm{fin}}$}
\]
Consequently the induced map
$\operatorname{coker}\kappa_0\to H^1(S_j,\mathcal O_{S_j})$ is an
isomorphism at each finite fibre.  Thus each summand
$T_{p_j}\to H^1(S_j,\mathcal O_{S_j})$ is an isomorphism.  In the long
sequence of (11.12), (11.16) first removes the entire degree-zero residue,
and (11.17) then cancels $T$.  What remains is precisely (11.13).
\end{proof}

\begin{remark}[the cusp is not the torsion locus]
At $p_0$, if
$Q=\nu_*\mathcal O_{\mathrm{dP}_6}/\mathcal O_W$, then
\[
 0\longrightarrow\Omega^1_{N_0/\Delta}\longrightarrow
 \mathcal E|_{N_0}\longrightarrow i_*Q\longrightarrow0,
 \qquad h^\bullet(Q)=(2,1,0).
\tag{11.17a}
\]
This relative sequence is the right column of the following exact
normalization--residue diagram; all sheaves in the bottom row are regarded as
sheaves on $N_0$ by $i_*$:
\[
\begin{array}{ccccccccc}
&&0&&0&&0&&\\
&&\big\downarrow&&\big\downarrow&&\big\downarrow&&\\
0&\longrightarrow&f^*\Omega^1_\Delta&\longrightarrow&\Omega^1_{N_0}
 &\longrightarrow&\Omega^1_{N_0/\Delta}&\longrightarrow&0\\
&&\big\downarrow&&\big\downarrow&&\big\downarrow&&\\
0&\longrightarrow&f^*\Omega^1_\Delta(\log0)&\longrightarrow&
 \Omega^1_{N_0}(\log W)&\longrightarrow&\mathcal E|_{N_0}
 &\longrightarrow&0\\
&&\big\downarrow\,{\scriptstyle\operatorname{res}_B}
&&\big\downarrow\,{\scriptstyle\operatorname{res}_W}
&&\big\downarrow\,{\scriptstyle\rho}&&\\
0&\longrightarrow&i_*\mathcal O_W&\longrightarrow&
 i_*\nu_*\mathcal O_{\mathrm{dP}_6}&\longrightarrow&i_*Q
 &\longrightarrow&0\\
&&\big\downarrow&&\big\downarrow&&\big\downarrow&&\\
&&0&&0&&0&&
\end{array}
\tag{11.17a$'$}
\]
The first two columns are the base and normalized-divisor residue sequences;
the bottom row is the definition of $Q$.  Hence the snake lemma identifies
the right column with (11.17a), including its displayed map $\rho$.

The exactness assertion is also visible stalkwise.  At a double chart $t=xy$, with one
coordinate along the double stratum and logarithmic generator $e$, the image
of ordinary relative differentials in $\mathcal E$ contains the stratum
differential and has transverse part $(x,y)e$; its cokernel is therefore
$R/(x,y)$.  At a triple chart $t=xyz$, write
$\mathcal E=Re_1\oplus Re_2$ with
$d\log x=e_1$, $d\log y=e_2$, and
$d\log z=-e_1-e_2$.  The image is generated by
\[
 xe_1,\qquad ye_2,\qquad z(e_1+e_2),
\]
Writing a normalized branch value as a triple modulo the diagonal, the
quotient map is explicitly
\[
 e_1\longmapsto[(1,0,0)],\qquad
 e_2\longmapsto[(0,1,0)].
\tag{11.17a$''$}
\]
Its kernel is exactly
$\langle xe_1,ye_2,z(e_1+e_2)\rangle$, so this is the standard local
presentation of the normalization quotient $Q$.
Away from the double locus the map is an isomorphism.  This proves (11.17a)
stalkwise.
Here are the maps, rather than only their ranks.  Quotienting the normalization
complex (9.10) by $\mathcal O_W$ gives
\[
 0\longrightarrow Q\longrightarrow
 \bigoplus_{k=1}^3\mathcal O_{D_k}
 \xrightarrow{\delta_1}\mathbb C_P\oplus\mathbb C_Q
 \longrightarrow0.
\]
Because $D_k\simeq\mathbb P^1$, its cohomology is the cohomology of
\[
 \mathbb C^3\longrightarrow\mathbb C^2,\qquad
 (a_1,a_2,a_3)\longmapsto(s,s),\qquad s=a_1-a_2+a_3.
\]
Thus
\[
 H^0(Q)=\{(a_1,a_2,a_3):a_1-a_2+a_3=0\}\simeq\mathbb C^2,
 \qquad H^1(Q)=\mathbb C^2/\mathbb C(1,1)\simeq\mathbb C,
\tag{11.17b}
\]
and $H^2(Q)=0$.  Since $\mathrm{dP}_6$ is rational, the connecting maps for
$0\to\mathcal O_W\to\nu_*\mathcal O_{\mathrm{dP}_6}\to Q\to0$ are
isomorphisms
\[
 \delta_0:H^0(Q)\xrightarrow{\sim}H^1(W,\mathcal O_W),
 \qquad
 \delta_1:H^1(Q)\xrightarrow{\sim}H^2(W,\mathcal O_W).
\tag{11.17c}
\]

At a triple chart $t=z_0z_1z_2$, a relative logarithmic form is represented
by $\eta=\sum a_i\,d\log z_i$, with holomorphic coefficient triple
$(a_0,a_1,a_2)$ taken modulo $\mathcal O(1,1,1)$.  Restricting the branch
differences to the three double strata gives
\[
 \rho(\eta)=(a_0-a_1,\ a_0-a_2,\ a_1-a_2).
\tag{11.17d}
\]
Formula (11.17d) is written in the $P$-ordering of the globally labelled
double curves.  The $Q$-point sign is fixed by a separate calculation, not by
an implicit relabelling.  In the common $e_1,e_2$-plane, use the ordered
branch vertices
\[
 (b_0,b_1,b_2)=(0,e_1,e_2),\qquad
 (c_0,c_1,c_2)=(e_1,e_1+e_2,e_2)
\]
at $P$ and $Q$, respectively.  The coefficient triples of
$d\log x_1,d\log x_2$ are
\[
 (0,1,0),\ (0,0,1)\quad\text{at }P,
 \qquad
 (1,1,0),\ (0,1,1)\quad\text{at }Q.
\]
The globally labelled branch-difference maps are
\[
 d_P(a)=(a_0-a_1,a_0-a_2,a_1-a_2),
 \qquad
 d_Q(a)=(a_2-a_1,a_0-a_1,a_0-a_2).
\tag{11.17d$'$}
\]
Direct substitution gives, at both triple points, the same two global
sections
\[
 q_1=(-1,0,1),\qquad q_2=(0,-1,-1).
\tag{11.17d$''$}
\]
For constant coefficients this induces an isomorphism
$\mathbb C^3/\mathbb C(1,1,1)\to H^0(Q)$.  In the toric chart
$(z_0,z_1,z_2)=(t/(x_1x_2),x_1,x_2)$ it sends
\[
 d\log x_1\longmapsto(-1,0,1),\qquad
 d\log x_2\longmapsto(0,-1,-1),
\]
which is a basis of the hyperplane in (11.17b); unimodular changes of cusp
chart merely change this basis by $\mathrm{GL}_2(\mathbb Z)$.

Write $\eta_1=d\log x_1$, $\eta_2=d\log x_2$, take the $q_i$ of
(11.17d$''$), and put $\alpha_i=\delta_0(q_i)$.  By (11.17c),
$\alpha_1,\alpha_2$ form a basis of $H^1(W,\mathcal O_W)$.  Naturality of
the connecting maps fixes a sign that must be retained.  With the standard
\v Cech differential
$(\delta c)_{i_0\ldots i_{r+1}}
=\sum_k(-1)^kc_{i_0\ldots\widehat{i_k}\ldots i_{r+1}}$, lift $q_i$ to the
normalization complex.  The total differential of a degree-one cocycle
$\alpha$ times that lift is
$(-1)^1\alpha\smile\delta\widetilde q_i$.  The projection formula therefore
gives exactly
\[
 \delta_1\bigl(\rho_*(\alpha\otimes\eta_i)\bigr)
 =-\alpha\smile\delta_0(q_i)=-\alpha\smile\alpha_i.
\]
The cup product in \cref{lem:cusp-structure-sheaf-audit} is an isomorphism.
Define, with this convention,
\[
 \theta=-\delta_1^{-1}(\alpha_1\smile\alpha_2).
\]
Graded commutativity now gives, with no residual sign ambiguity,
\[
 \rho_*=(0,1,-1,0):
 \mathbb C\langle
 \alpha_1\eta_1,\alpha_1\eta_2,\alpha_2\eta_1,\alpha_2\eta_2
 \rangle\longrightarrow\mathbb C\theta.
\tag{11.17e}
\]
In particular the map is surjective.  Applying direct images to (11.17a),
the first map is reduction modulo $t$ followed by the isomorphism (11.17d),
and the second is (11.17e).  Consequently
\[
 f_*\Omega^1_{N_0/\Delta}=t\mathcal O^2,\quad
 R^1f_*\Omega^1_{N_0/\Delta}=\ker(\mathcal O^4\to\mathbb C_{p_0}),
 \quad R^2f_*\Omega^1_{N_0/\Delta}=\mathcal O^2.
\]
For the row in (11.17e), the middle kernel is generated by
$e_1,e_4,e_2+e_3,te_2$.
Thus $\Omega^1_{X/B}$ is torsion-free in a neighbourhood of the cusp, but
non-locally-free along its $xy$ and $xyz$ strata.  Globally its torsion is
supported precisely on the two multiple fibres, as the following local module
calculation proves.  On the smooth cover $J'_j$, let
$R=\mathcal O_{J'_j,x}$ and choose coordinates $(s,y_1,y_2)$ smooth over the
$s$-disc.  Relative to the original base $B$, where $t=s^{m_j}$, one has
\[
 \Omega^1_{J'_j/B,x}
 =\frac{R\,ds\oplus R\,dy_1\oplus R\,dy_2}
        {(m_js^{m_j-1}ds)}
 \simeq R\,dy_1\oplus R\,dy_2
 \oplus\bigl(R/(s^{m_j-1})\bigr)ds.
\tag{11.17f}
\]
Thus the last summand is exactly the torsion submodule.  The quotient map
$q:J'_j\to N_j$ is finite \'{e}tale, so the relative cotangent sequence gives
\[
 q^*\Omega^1_{N_j/B}\simeq\Omega^1_{J'_j/B}.
\]
Faithful flatness descends both the presence of the displayed torsion and its
absence off $s=0$.  Hence ordinary relative torsion is supported exactly on
the two multiple fibres and is absent at the cusp.  The reconstructed
manuscript's assertion that torsion occurs at all three singular fibres is
false.
\end{remark}

\subsection{The scalar and the exact Hodge diamond}

\begin{theorem}[the remaining scalar is nonzero]
\label{thm:h12-scalar}
\Verified\ The induced map
\[
 H^0(B,R^1f_*\mathcal E)\simeq\mathbb C^2
 \xrightarrow{\ H^0(\kappa_1)\ }
 H^0(B,\mathcal O)\simeq\mathbb C
\]
has rank one.
\end{theorem}

\begin{proof}
Every matrix $R_g$ in (2.3d) has second column $(0,1)^{\mathsf t}$.
Consequently the marked vertical field
\[
 \xi=\frac{\partial}{\partial\zeta_2}
\]
is invariant under every deck transformation.  The affine translations in
\cref{prop:candidate-finite-fillings} have identity differential in this
direction, and the corresponding toric logarithmic direction extends across
the cusp charts of \cref{thm:candidate-cusp-geometry}.  Thus $\xi$ is a
global section nowhere zero in the locally free relative sheaf
$T_{X/B}=\mathcal E^\vee$ and is the section used in the evaluation sequence
\[
 0\longrightarrow\mathcal O(-1)\longrightarrow f_*\mathcal E
 \xrightarrow{\operatorname{ev}_\xi}\mathcal O\longrightarrow0.
\]
This sequence splits, and the lift $h$ of $1\in H^0(\mathcal O)$ is unique:
the difference of two lifts is a global section of $\mathcal O(-1)$, while
$H^0(B,\mathcal O(-1))=0$.  On the marked cover write that unique lift as
\[
 h=(a(z),1)\binom{d\zeta_1}{d\zeta_2}.
\tag{11.18}
\]
The function $a$ is holomorphic on each finite cover because $h$ is a
section of the actual invariant lattice, not merely of its generic fibre.
For $g_2$, invariance of (11.18) and
$R_{g_2}=\left(\begin{smallmatrix}1/\tau&0\\-\mu/\tau&1\end{smallmatrix}\right)$
mean exactly
\[
 (a(z),1)=(a(g_2z),1)R_{g_2}(z),
 \qquad\text{hence}\qquad
 a(g_2z)=\tau(z)a(z)+\mu(z).
\tag{11.18a}
\]
At the fixed point $z_2$, the normalized values
$\tau(z_2)=i$ and $\mu(z_2)=(1-i)/2$ therefore give
\[
 a(z_2)=\frac12.
\tag{11.18b}
\]
The cup-product map
$f_*\mathcal E\otimes G\to R^1f_*\mathcal E$ sends
$h\otimes\gamma$ to the unit generator $e_0b_0$ of the degree-zero
character line in (11.4).  This identifies the precise entry of the row
(11.11) that must be evaluated.

Normalize the order-four linearizing coordinate by
\[
 s=(z_2-\bar z_2)s_2,
 \qquad s\circ g_2=-is,
 \qquad \frac{ds}{dz}(z_2)=1.
\]
Write a dot for $d/ds$, and put
$\ell=dt/t=4\,ds/s$.  In the quotient basis $(\gamma,u)$ of $G$, the
marked period derivative and $ds=(s/4)\ell$ give the exact logarithmic map
\[
 \kappa_0(h)=\frac{s}{4}
 \Bigl[(6a\dot\mu+\dot\beta)\gamma
 +(a\dot\tau+\dot\mu)u\Bigr]\otimes\ell.
\tag{11.19}
\]
For $(m,a)=(4,1)$, the actual invariant target generator from (11.4) is
$s(b_0\wedge b_1)\ell$.  Normalize $b_1$ so that
\[
 \gamma\wedge b_1=u_2(s)(\gamma\wedge u),
 \qquad u_2(0)=1;
\]
equivalently, use the invariant frame
\[
 \omega_2=s\,u_2(s)(\gamma\wedge u)\ell.
\tag{11.19f}
\]
The uncorrected expression $s(\gamma\wedge u)\ell$ is only invariant at
$s=0$: indeed $g_2^*(\gamma\wedge u)=\tau(z)(\gamma\wedge u)$, so its
pullback factor is $(-i)\tau(z)$.  Formula (11.11), in the genuine invariant
frame (11.19f), sends the first ordered source generator $e_0b_0$ to
minus the second coefficient of $\kappa_0(e_0)$.  Since $u_2(0)=1$,
(11.19) therefore gives the exact signed coarse value
\[
 \kappa_1(h\otimes\gamma)(p_2)
 =-\frac14\bigl(a(z_2)\dot\tau(z_2)+\dot\mu(z_2)\bigr).
\tag{11.19a}
\]

At the order-four lift $z_2$, the quotient coordinate has order four while
the modular invariant has order two at $\tau=i$; see
\cite[Chapter~VII]{Serre1973}.  Equivalently,
$j(\tau)-1728$ has a double zero at $i$ and
$j(\tau(z))-1728$ has order four at $z_2$.  Hence
\[
 \operatorname{ord}_{z_2}(\tau-i)=2,\qquad \dot\tau(z_2)=0.
\tag{11.20}
\]
By (11.7), $K(z_2)\ne0$; since
$K=6\dot\mu^2-\dot\tau\dot\beta$, (11.20) gives
$K(z_2)=6\dot\mu(z_2)^2\ne0$.  Substitution in (11.19a) therefore gives
\[
 \kappa_1(h\otimes\gamma)(p_2)=-\frac14\dot\mu(z_2)\ne0.
\tag{11.20a}
\]
The image is a global section of
$\mathcal L\otimes R^2f_*\mathcal O_X\simeq\mathcal O_B$; because it is
nonzero at $p_2$, it is a nonzero constant.  Hence
$H^0(\kappa_1)$ has rank one.
\end{proof}

\begin{corollary}[exact Hodge diamond]
\label{cor:exact-hodge-diamond}
\Verified\ The last parameter in
\cref{thm:parametric-hodge-diamond} is
\[
 h^{1,2}(X)=1.
\]
Consequently $h^{1,1}=2$, and the Hodge table is
\[
\begin{array}{c|cccc}
q\backslash p&0&1&2&3\\ \hline
3&0&0&0&1\\
2&0&1&2&1\\
1&1&2&1&0\\
0&1&0&0&0.
\end{array}
\]
The two nontrivial Fr\"olicher $d_1$ rows are both exact complexes
$0\to\mathbb C\to\mathbb C^2\to\mathbb C\to0$, and
$\dim E_1=10$, $\dim E_2=2$.
\end{corollary}

\begin{proof}
Put $\mathcal K=\ker\kappa_1$.  Since
\[
0\longrightarrow\mathcal K\longrightarrow
\mathcal O^{\oplus2}\oplus\mathcal O(-1)^{\oplus2}
\xrightarrow{\kappa_1}\mathcal O\longrightarrow0
\]
and the global map has rank one, its restriction to the two trivial summands
is a nonzero constant row (the negative summands have no global sections).
After a constant change of basis it is the identity on one copy of
$\mathcal O$.  Subtracting the remaining component from that summand makes
the map the projection, so its kernel is
\[
 \mathcal K\simeq\mathcal O\oplus\mathcal O(-1)^{\oplus2}.
\]
By 
\cref{prop:absolute-cotangent-direct-images} and coherent Leray,
\[
 h^{1,2}=h^1(B,\mathcal K)
 +h^0(B,\mathcal O\oplus\mathcal O(-1))=0+1=1.
\]
The remaining assertions follow from
\cref{thm:parametric-hodge-diamond,thm:frolicher-e2-candidate}.
\end{proof}

\subsection{Bott--Chern and Aeppli dimensions}

For a compact complex manifold write
\[
 H^{p,q}_{BC}(X)=
 \frac{\ker\partial\cap\ker\bar\partial\subset A^{p,q}(X)}
      {\operatorname{im}\partial\bar\partial},
 \qquad
 H^{p,q}_{A}(X)=
 \frac{\ker\partial\bar\partial\subset A^{p,q}(X)}
      {\operatorname{im}\partial+\operatorname{im}\bar\partial}.
\]

\begin{theorem}[exact Bott--Chern and Aeppli diamonds]
\label{thm:exact-bott-chern-aeppli}
\Verified\ The Bott--Chern dimensions of the candidate are
\[
\begin{array}{c|cccc}
q\backslash p&0&1&2&3\\ \hline
3&0&0&1&1\\
2&0&1&2&1\\
1&0&2&1&0\\
0&1&0&0&0,
\end{array}
\tag{11.21}
\]
and its Aeppli dimensions are
\[
\begin{array}{c|cccc}
q\backslash p&0&1&2&3\\ \hline
3&0&0&0&1\\
2&0&1&2&0\\
1&1&2&1&0\\
0&1&1&0&0.
\end{array}
\tag{11.22}
\]
In total degree $k=0,\ldots,6$ this is
\[
 (h^k_{BC})=(1,0,2,2,2,2,1),\qquad
 (h^k_A)=(1,2,2,2,2,0,1).
\tag{11.23}
\]
Consequently, set
\[
 \Delta^k=h^k_{BC}+h^k_A-2b_k
\]
for the degree-$k$ non-$\partial\bar\partial$ defects.  The
Angella--Tomassini inequality
\cite[Theorem~A]{AngellaTomassini2013} gives $\Delta^k\ge0$ for every compact
complex manifold; here the exact values are
\[
 (\Delta^k)=(0,2,4,4,4,2,0).
\tag{11.24}
\]
\end{theorem}

\begin{proof}
McHugh associates to a compact complex manifold the finite-dimensional
quotient
\[
 K^{p,q}=
 \frac{\ker\partial\cap\ker\bar\partial\subset A^{p,q}}
 {\bigl(\operatorname{im}\partial\cap\ker\bar\partial\bigr)
  +\bigl(\operatorname{im}\bar\partial\cap\ker\partial\bigr)},
 \qquad k^{p,q}=\dim_{\mathbb C}K^{p,q}.
\]
Put
\[
 Z^{p,q}=\ker\partial\cap\ker\bar\partial\subset A^{p,q},
 \qquad
 B_K^{p,q}=
 \bigl(\operatorname{im}\partial\cap\ker\bar\partial\bigr)
 +\bigl(\operatorname{im}\bar\partial\cap\ker\partial\bigr).
\]
The natural map $Z^{p,q}\to H^{p+q}_{\mathrm{dR}}(X;\mathbb C)$ need not
descend to an injection from $K^{p,q}$; the dimension estimate has the
opposite quotient mechanism.  If a pure form $\alpha\in Z^{p,q}$ is
$d$-exact, the $(p,q)$ component of a primitive gives
\[
 \alpha=\partial\beta^{p-1,q}+\bar\partial\beta^{p,q-1}.
\]
Because $\bar\partial\alpha=0$ and $\partial\alpha=0$, respectively,
$\partial\beta^{p-1,q}$ is $\bar\partial$-closed and
$\bar\partial\beta^{p,q-1}$ is $\partial$-closed.  Thus
$\ker(Z^{p,q}\to H^{p+q}_{\mathrm{dR}})\subset B_K^{p,q}$, so the de Rham
image of $Z^{p,q}$ surjects onto $K^{p,q}$.  Hence
\[
 k^{p,q}\le b_{p+q}.
\tag{11.25}
\]
Because $X$ has the integral cohomology of $S^6$, (11.25) gives
$k^{p,q}=0$ whenever $0<p+q<6$.

For a compact complex threefold, McHugh's Theorem~6.3,
Proposition~6.4, equation~(7), and the ensuing table give
\begin{align*}
 h^{1,1}_{BC}
 &=2h^{0,1}-b_1+\dim_{\mathbb R}H^{1,1}(X;\mathbb R),\\
 h^{1,2}_{BC}
 &=h^{1,2}+\dim_{\mathbb C}E_2^{3,1}-k^{2,0},\\
 h^{2,2}_{BC}
 &=-h^{1,1}_{BC}+h^{0,1}-h^{0,2}-h^{1,0}+h^{1,1}
   +h^{1,2}+h^{2,0}\\
 &\hspace{2em}+2\dim_{\mathbb C}E_2^{3,1}
   +k^{1,1}-k^{1,2}-2k^{2,0},
\end{align*}
and his threefold table gives the remaining entries
\cite[Theorem~6.3, Proposition~6.4, equation~(7), \S7.1]{McHugh2019}; see also the
independent synthesis in \cite[\S4]{Angella2018}.  Here
$H^{1,1}(X;\mathbb R)$ denotes the subspace of
$H^2_{\mathrm{dR}}(X;\mathbb R)$ represented by real $d$-closed forms of
type $(1,1)$, and it is zero because $H^2_{\mathrm{dR}}(X;\mathbb R)=0$,
while \cref{thm:frolicher-e2-candidate} gives $E_2=E_\infty$; since
$b_4=0$, this in particular gives $E_2^{3,1}=0$.  Substitution of
\[
 h^{0,1}=1,\quad h^{0,2}=h^{1,0}=h^{2,0}=0,\quad
 h^{1,1}=2,\quad h^{1,2}=1
\]
therefore gives
\[
 h^{1,1}_{BC}=2,\qquad h^{1,2}_{BC}=h^{2,1}_{BC}=1,\qquad
 h^{2,2}_{BC}=2.
\]
The same table gives
$h^{2,3}_{BC}=h^{3,2}_{BC}=1$, the two endpoint dimensions equal to one,
and every other entry zero.  This proves (11.21).  The Hodge-star duality
of Schweitzer \cite[\S2.c, pp.~9--10]{Schweitzer2007},
$h_A^{p,q}=h_{BC}^{3-p,3-q}$, proves (11.22).  Summing along diagonals and
using $b_0=b_6=1$ and $b_k=0$ for $1\le k\le5$ proves
(11.23)--(11.24).
\end{proof}

The dimension calculation above is completed by the explicit representative
and product theorem in the following section.

\section{Period-coordinate Dolbeault classes and the Bott--Chern algebra}
\label{sec:bcpr-product-structure}

\subsection{Fixed input and conventions}

Let $f\colon X\to B=\mathbb P^1$ be the compact complex threefold of
\cref{thm:candidate-global-geometry}.  We use
$d=\partial+\bar\partial$ and, following
\cite[\S1.a]{Schweitzer2007},
\begin{equation}
 H^{p,q}_{BC}(X)=
 \frac{\ker\partial\cap\ker\bar\partial
       \subset\mathcal A^{p,q}(X)}
      {\operatorname{im}(\partial\bar\partial)}.
\label{eq:bcpr-bc-definition}
\end{equation}
The natural map to Dolbeault cohomology is denoted
\begin{equation}
 \pi_{\bar\partial}^{p,q}\colon H^{p,q}_{BC}(X)
 \longrightarrow H^{p,q}_{\bar\partial}(X).
\label{eq:bcpr-dolbeault-map}
\end{equation}
Exterior multiplication makes the direct sum of the Bott--Chern groups a
bigraded algebra, and complex conjugation interchanges bidegrees $(p,q)$
and $(q,p)$ \cite[\S1.a]{Schweitzer2007}.

By \cref{cor:exact-hodge-diamond}, the two nonzero rows of the
Fr\"olicher differential are the exact complexes
\begin{equation}
 0\longrightarrow H^{0,1}_{\bar\partial}(X)
 \xrightarrow{d_1}H^{1,1}_{\bar\partial}(X)
 \xrightarrow{d_1}H^{2,1}_{\bar\partial}(X)
 \longrightarrow0,
\label{eq:bcpr-hodge-row-one}
\end{equation}
\begin{equation}
 0\longrightarrow H^{1,2}_{\bar\partial}(X)
 \xrightarrow{d_1}H^{2,2}_{\bar\partial}(X)
 \xrightarrow{d_1}H^{3,2}_{\bar\partial}(X)
 \longrightarrow0,
\label{eq:bcpr-hodge-row-two}
\end{equation}
and each row has dimensions $1,2,1$.  The exact Bott--Chern calculation
of \cref{thm:exact-bott-chern-aeppli}, obtained from the compact-threefold
formulae of \cite{McHugh2019}, gives
\begin{equation}
\begin{array}{c|cccccccc}
(p,q)&(0,0)&(1,1)&(2,1)&(1,2)&(2,2)&(3,2)&(2,3)&(3,3)\\ \hline
h^{p,q}_{BC}&1&2&1&1&2&1&1&1.
\end{array}
\label{eq:bcpr-nonzero-dimensions}
\end{equation}
All other Bott--Chern groups vanish.  By
\cref{cor:candidate-sphere}, $b_k(X)=0$ for $1\leq k\leq5$.

On a marked smooth fibre retain the integral basis
\[
 V=\mathbb Z\langle\gamma,u,w,\delta\rangle
\]
without renaming its element $u$, and put
\begin{equation}
 \eta=u\wedge w+6\gamma\wedge\delta,
 \qquad
 \eta^2=12\gamma\wedge u\wedge w\wedge\delta.
\label{eq:bcpr-eta-square}
\end{equation}
The holomorphic Hodge rows supplied by
\cref{prop:period-input-audit,prop:very-general-torus-audit} are
\begin{equation}
 \sigma_1=6\mu\gamma+\tau u+w,
 \qquad
 \sigma_2=\beta\gamma+\mu u+\delta.
\label{eq:bcpr-period-rows}
\end{equation}
Thus no global Dolbeault or Bott--Chern generator below is denoted by
$u$ or $x$; those letters remain reserved for the printed fibre basis and
for local real linear functions.

\subsection{Exact fibre representatives}

Write
\begin{equation}
 y=\operatorname{Im}\tau,
 \qquad m=\operatorname{Im}\mu,
 \qquad r=\operatorname{Im}\beta,
 \qquad \Delta=r-\frac{6m^2}{y}<0.
\label{eq:bcpr-delta}
\end{equation}
The four marked lattice columns in the fibre coordinate
$\zeta=(\zeta_1,\zeta_2)$ are
\begin{equation}
 (6\mu,\beta),\quad(\tau,\mu),\quad(1,0),\quad(0,1),
\label{eq:bcpr-lattice-columns}
\end{equation}
in the order $\widehat\gamma,\widehat u,\widehat w,
\widehat\delta$.  On the universal cover of a marked fibre define the
real linear functions
\begin{equation}
 x_\gamma=
 \frac{y\,\operatorname{Im}\zeta_2-m\,\operatorname{Im}\zeta_1}
      {y\Delta},
 \qquad
 x_u=
 \frac{r\,\operatorname{Im}\zeta_1-6m\,\operatorname{Im}\zeta_2}
      {y\Delta}.
\label{eq:bcpr-real-linear-functions}
\end{equation}
For the first two columns in \cref{eq:bcpr-lattice-columns}, the numerators
of the changes in $x_\gamma$ are respectively
$yr-6m^2=y\Delta$ and $ym-my=0$, while the numerators of the changes in
$x_u$ are $6rm-6mr=0$ and $ry-6m^2=y\Delta$.  The last two columns are
real and change neither imaginary part.  Hence
\[
 \bigl(\Delta_{\widehat\gamma}x_\gamma,
       \Delta_{\widehat u}x_\gamma,
       \Delta_{\widehat w}x_\gamma,
       \Delta_{\widehat\delta}x_\gamma\bigr)=(1,0,0,0),
\]
\begin{equation}
 \bigl(\Delta_{\widehat\gamma}x_u,
       \Delta_{\widehat u}x_u,
       \Delta_{\widehat w}x_u,
       \Delta_{\widehat\delta}x_u\bigr)=(0,1,0,0).
\label{eq:bcpr-linear-function-increments}
\end{equation}
Their differentials, although not the functions themselves, therefore
descend to the fibre.  Their $(0,1)$ parts are
\begin{equation}
 \Gamma:=\bar\partial_Fx_\gamma
 =\frac{i}{2\Delta}
 \left(d\bar\zeta_2-\frac{m}{y}d\bar\zeta_1\right),
\label{eq:bcpr-gamma-form}
\end{equation}
\begin{equation}
 U:=\bar\partial_Fx_u
 =\frac{i}{2y\Delta}
 \left(r\,d\bar\zeta_1-6m\,d\bar\zeta_2\right).
\label{eq:bcpr-u-form}
\end{equation}
They represent respectively the $(0,1)$ images of the integral classes
$\gamma$ and $u$.

The identity
\begin{equation}
 \eta=u\wedge\sigma_1+6\gamma\wedge\sigma_2
     =-\sigma_1\wedge u-6\sigma_2\wedge\gamma
\label{eq:bcpr-eta-decomposition}
\end{equation}
holds in $\bigwedge^2V_{\mathbb C}$.  Since $\eta$ is of type $(1,1)$
on every marked fibre, \cref{eq:bcpr-eta-decomposition} gives its exact
Dolbeault representative
\begin{equation}
 Q_F=-d\zeta_1\wedge U-6d\zeta_2\wedge\Gamma.
\label{eq:bcpr-q-fibre}
\end{equation}
Equivalently, with the Appell--Humbert convention
$\eta=(2i)^{-1}\sum(H_\eta)_{ab}\sigma_a\wedge\bar\sigma_b$, the matrix
calculation in the proof of \cref{prop:very-general-torus-audit} gives
\begin{equation}
 H_\eta=\frac1{y\Delta}
 \begin{pmatrix}r&-6m\\-6m&6y\end{pmatrix},
\label{eq:bcpr-appell-humbert-matrix}
\end{equation}
and therefore
\begin{equation}
 Q_F=\frac1{2iy\Delta}\Bigl(
 r\,d\zeta_1\wedge d\bar\zeta_1
 -6m\,d\zeta_1\wedge d\bar\zeta_2
 -6m\,d\zeta_2\wedge d\bar\zeta_1
 +6y\,d\zeta_2\wedge d\bar\zeta_2\Bigr).
\label{eq:bcpr-q-fibre-expanded}
\end{equation}
Direct exterior multiplication, with holomorphic differentials ordered
before antiholomorphic differentials, gives
\begin{equation}
 Q_F\wedge\Gamma
 =\frac1{4y\Delta}
 d\zeta_1\wedge d\bar\zeta_1\wedge d\bar\zeta_2\neq0,
\label{eq:bcpr-q-gamma-product}
\end{equation}
\begin{equation}
 Q_F^2
 =\frac3{y\Delta}
 d\zeta_1\wedge d\zeta_2\wedge
 d\bar\zeta_1\wedge d\bar\zeta_2\neq0.
\label{eq:bcpr-q-square}
\end{equation}
Indeed,
$\det H_\eta=(6ry-36m^2)/(y^2\Delta^2)=6/(y\Delta)$,
and the coefficient in \cref{eq:bcpr-q-square} is
$\tfrac12\det H_\eta$ in the displayed ordering.  The identities
\Cref{eq:bcpr-q-gamma-product,eq:bcpr-q-square} will be used only in
fibre cohomology.

\subsection{Extension of the two Dolbeault inputs}

\begin{lemma}[the degree-one input]
\label{lem:bcpr-degree-one-input}
\Verified\ There is a distinguished class
$\mathfrak g\in H^{0,1}_{\bar\partial}(X)$ whose restriction to every
smooth marked fibre is represented by $\Gamma$.
\end{lemma}

\begin{proof}
By \cref{thm:structure-direct-images-audit}, the monodromy-invariant
integral class $\gamma$ maps to the nowhere-zero holomorphic section
$\gamma^{0,1}$ of
$R^1f_*\mathcal O_X\simeq\mathcal O\oplus\mathcal O(-1)$.  Coherent
Leray gives
\[
 H^1(X,\mathcal O_X)=H^0(B,R^1f_*\mathcal O_X)
 =\mathbb C\gamma^{0,1}.
\]
Let $\mathfrak g$ be the class mapping to $\gamma^{0,1}$.
\Cref{eq:bcpr-gamma-form} is the exact period-coordinate
representative of its fibre restriction.  This is the distinguished
coordinate of \cref{cor:picard-line-audit} before any analytic rescaling.
\end{proof}

\begin{lemma}[extension of the invariant $(1,1)$ class]
\label{lem:bcpr-eta-extension}
\Verified\ The fibre class $\eta$ extends to a nonzero holomorphic section
\begin{equation}
 \boldsymbol\eta\in
 H^0\bigl(B,R^1f_*\Omega_X^1\bigr).
\label{eq:bcpr-eta-global-section}
\end{equation}
Consequently there is a class
$\mathfrak q\in H^{1,1}_{\bar\partial}(X)$ whose coherent-Leray edge
value is $\boldsymbol\eta$.  For every smooth fibre inclusion
$i_b\colon F_b\to X$,
\begin{equation}
 i_b^*\mathfrak q=[Q_F]_{\bar\partial}
 \in H^{1,1}_{\bar\partial}(F_b).
\label{eq:bcpr-q-restriction}
\end{equation}
\end{lemma}

\begin{proof}
Over the smooth marked locus, $\eta$ is flat, monodromy invariant, and of
type $(1,1)$ identically.  Equation
\Cref{eq:bcpr-eta-decomposition} identifies the corresponding holomorphic
section of the relative group as
\begin{equation}
 -\sigma_1\wedge u^{0,1}-6\sigma_2\wedge\gamma^{0,1}.
\label{eq:bcpr-relative-eta-section}
\end{equation}

At a finite filling write $t=s^M$, with $(M,a)=(3,2)$ or $(4,1)$.
The covering family is smooth at $s=0$, and
\cref{eq:bcpr-eta-decomposition} holds identically there, so
\cref{eq:bcpr-relative-eta-section} remains holomorphic.  It is invariant
under the cyclic action.  In equivariant frames $e_0,e_1$ for the relative
cotangent and $b_0,b_1$ for $R^1\mathcal O$, the actual invariant
$\mathbb C\{t\}$-basis from
\cref{prop:log-relative-direct-images} is
\begin{equation}
 e_0b_0,\quad s^ae_0b_1,\quad
 s^{M-a}e_1b_0,\quad e_1b_1.
\label{eq:bcpr-finite-invariant-basis}
\end{equation}
Exactness of cyclic invariants in characteristic zero therefore makes
\cref{eq:bcpr-relative-eta-section} a holomorphic descended section at
both finite points.

At the cusp, $\gamma$ and $u$ extend as sections of $R^1f_*\mathbb C$,
and their $(0,1)$ images form a basis of $H^1(W,\mathcal O_W)$ by
\cref{thm:structure-direct-images-audit}.  On the dense torus,
\[
 \sigma_i=d\zeta_i=(2\pi i)^{-1}d\log x_i.
\]
The classes of these logarithmic forms extend in the relative logarithmic
sheaf as
\begin{equation}
 \mathcal E|_{N_0}
 =\mathcal O_{N_0}[d\log x_1]\oplus
  \mathcal O_{N_0}[d\log x_2]
\label{eq:bcpr-cusp-log-basis}
\end{equation}
by \cref{prop:log-relative-direct-images}.  A cusp quotient shear changes
$d\log x_i$ only by a holomorphic multiple of $d\log t$, which is zero in
the relative logarithmic sheaf.  Thus
\cref{eq:bcpr-relative-eta-section} also extends holomorphically through
the cusp.  This concerns classes in the relative logarithmic sheaf; it
does not assert that the $d\log x_i$ are global ordinary one-forms on $X$.

The Gauss--Manin/Kodaira--Spencer map
\[
 \kappa_1\colon R^1f_*\mathcal E\longrightarrow
 \Omega_B^1(\log\Sigma)\otimes R^2f_*\mathcal O_X
\]
annihilates \cref{eq:bcpr-relative-eta-section} on the smooth locus: a
flat class which remains of type $(1,1)$ has zero $(0,2)$ Higgs component.
Both the section and $\kappa_1$ extend holomorphically, so the identity
theorem gives $\kappa_1(\boldsymbol\eta)=0$ on $B$.  The ordinary-cotangent
calculation in
\cref{prop:log-ks-divisor,prop:absolute-cotangent-direct-images,thm:h12-scalar}
identifies
\begin{equation}
 R^1f_*\Omega_X^1=\ker\kappa_1=:\mathcal K,
 \qquad
 \mathcal K\simeq\mathcal O\oplus\mathcal O(-1)^{\oplus2}.
\label{eq:bcpr-cotangent-kernel}
\end{equation}
This proves \cref{eq:bcpr-eta-global-section}.  Its value on every smooth
fibre is nonzero in Dolbeault cohomology, not merely as a differential form:
by \cref{eq:bcpr-eta-square} its square is the integral top class
$12\gamma\wedge u\wedge w\wedge\delta$, and therefore
\begin{equation}
 \int_{(F_b,\mathrm{or}_{\rm mark})}Q_F^2=12,
 \qquad
 \int_{(F_b,\mathrm{or}_{\mathbb C})}Q_F^2=-12
\qquad(b\in B^\circ),
\label{eq:bcpr-q-square-integral}
\end{equation}
where $\mathrm{or}_{\rm mark}$ is the orientation determined by
$(\gamma,u,w,\delta)$ and $\mathrm{or}_{\mathbb C}$ is the complex
orientation.  They are opposite because the determinant of the four marked
lattice columns in real coordinates is $y\Delta<0$.  Thus
$[Q_F]_{\bar\partial}\ne0$, so the holomorphic
section $\boldsymbol\eta$ is nonzero.  Explicitly, if
$Q_F=\bar\partial_F\alpha$ for a $(1,0)$-form $\alpha$, then
$Q_F^2=\bar\partial_F(\alpha\wedge Q_F)$; integration and Stokes' theorem
would contradict \cref{eq:bcpr-q-square-integral}.  Here
$\alpha\wedge Q_F$ has bidegree $(2,1)$, so its $\partial_F$ is zero in
complex dimension two and the displayed $\bar\partial_F$-boundary is also a
$d$-boundary.

Finally, coherent Leray over the curve gives the exact edge sequence
\begin{equation}
 0\longrightarrow H^1(B,f_*\Omega_X^1)
 \longrightarrow H^1(X,\Omega_X^1)
 \longrightarrow H^0(B,R^1f_*\Omega_X^1)
 \longrightarrow0,
\label{eq:bcpr-cotangent-edge}
\end{equation}
where $f_*\Omega_X^1\simeq\mathcal O(-2)$ by
\cref{prop:absolute-cotangent-direct-images}.  Choose $\mathfrak q$ over
$\boldsymbol\eta$.  Base change on a smooth fibre gives
\cref{eq:bcpr-q-restriction}.  Any two choices differ by the
one-dimensional left term in \cref{eq:bcpr-cotangent-edge}.  That term
restricts to zero on every smooth fibre.  The nonzero class
$d_1\mathfrak g$ also restricts to zero, because its restriction is
$[\partial_F\Gamma]_{\bar\partial}=0$; since
$H^1(B,\mathcal O(-2))$ is one-dimensional, the left term is exactly
$\operatorname{im}(d_1\colon H^{0,1}_{\bar\partial}\to
H^{1,1}_{\bar\partial})$.
\end{proof}

\begin{remark}[what is and is not explicit]
\label{rem:bcpr-literal-gauge}
Choose global $\bar\partial$-closed forms
\begin{equation}
 G\in\mathcal A^{0,1}(X),\qquad
 Q\in\mathcal A^{1,1}(X)
\label{eq:bcpr-global-representatives}
\end{equation}
representing $\mathfrak g$ and $\mathfrak q$.
\Cref{eq:bcpr-q-restriction} says that
$i_b^*[Q]_{\bar\partial}$ is the Dolbeault class of the explicit
translation-invariant
form displayed in \cref{eq:bcpr-q-fibre-expanded}; it does not say that an arbitrary
representative $Q$ restricts literally to that form.  A literal global
coefficient formula would require additional $C^\infty$ gauge choices and
horizontal correction forms on the disc--pair-of-pants cover, compatible
with the two affine cyclic conjugacies and the toric cusp.  Those choices
are not part of the canonical data.  No such literal gauge is needed for
the class and product calculations below.
\end{remark}

\subsection{The factor \texorpdfstring{$12$}{12} and the Bott--Chern first Chern class}

Let
\[
 \exp\colon H^1(X,\mathcal O_X)\longrightarrow\operatorname{Pic}(X)
\]
be induced by $z\mapsto e^{2\pi iz}$.  By
\cref{prop:leray-e2},
\begin{equation}
 H^0(B,R^1f_*\mathbb Z)=\mathbb Z(12\gamma).
\label{eq:bcpr-global-integral-line}
\end{equation}
The exponential--Leray comparison gives a line bundle $M$ on $B$ such
that
\begin{equation}
 \exp(12\mathfrak g)=f^*M,
 \qquad c_1(M)=d_2(12\gamma).
\label{eq:bcpr-exponential-transgression}
\end{equation}
For completeness, the first assertion follows because
\cref{eq:bcpr-global-integral-line} says that the restriction of
$12\mathfrak g$ to every fibre lies in the image of integral cohomology
and hence exponentiates to the trivial bundle.  The low-degree Leray
sequence for $\mathcal O^*$ then makes the resulting bundle a pullback.
On a disc--pair-of-pants cover of $B$, choose local integral lifts of
$12\gamma$.  Their differences on the annular overlaps form a Cech
one-cochain, and its Cech coboundary is simultaneously the two-cocycle
defining the Leray transgression and the two-cocycle of logarithms defining
$c_1(M)$.  This proves the second assertion with the same sign convention;
it is the exponential-sequence version of the local-lift calculation in
the proof of \cref{lem:first-transgression}.

Orient the generator $\omega\in H^2(B;\mathbb Z)$ by
\begin{equation}
 \omega=c_1(\mathcal O_B(1)).
\label{eq:bcpr-base-orientation}
\end{equation}
Then \cref{lem:first-transgression} and
\cref{eq:bcpr-exponential-transgression} give
\begin{equation}
 \exp(12\mathfrak g)=f^*\mathcal O_B(1).
\label{eq:bcpr-factor-twelve}
\end{equation}
If the opposite orientation is used for the primitive generator in
\cref{lem:first-transgression}, the corresponding statement is
$\exp(-12\mathfrak g)=f^*\mathcal O_B(1)$.  Thus the canonical integral
data fixes the absolute factor $12$, while its sign is the orientation
choice fixed in \cref{eq:bcpr-base-orientation}.

Put
\begin{equation}
 a=\partial G,
 \qquad a^\dagger=\bar\partial\bar G,
 \qquad
 \mathsf A=[a]_{BC},
 \qquad \mathsf B=[a^\dagger]_{BC}.
\label{eq:bcpr-a-b-definitions}
\end{equation}
Both forms are $d$-closed.  We now compute their Chern-class normalization
without choosing local metric weights.  For $z\in\mathbb C$, equip the
smooth trivial line with a global frame $s$ and the holomorphic structure
\begin{equation}
 \bar\partial_zs=-2\pi i zG\otimes s.
\label{eq:bcpr-holomorphic-structure}
\end{equation}
It is integrable because $\bar\partial G=0$.  If
$G=\bar\partial\phi_i$ on a good cover, then
$e_i=e^{2\pi iz\phi_i}s$ is $\bar\partial_z$-holomorphic and
\[
 e_j=e^{2\pi iz(\phi_j-\phi_i)}e_i.
\]
The holomorphic cocycle $\phi_j-\phi_i$ represents $\mathfrak g$ under
the Cech--Dolbeault comparison, so
\cref{eq:bcpr-holomorphic-structure} is precisely the line bundle
$\exp(z\mathfrak g)$ for the displayed exponential convention.

Give $s$ unit length.  The unit-metric Chern connection is
\begin{equation}
 \nabla_z=d-2\pi i\bigl(zG+\bar z\,\bar G\bigr).
\label{eq:bcpr-unit-chern-connection}
\end{equation}
Indeed, its $(0,1)$ part is
\cref{eq:bcpr-holomorphic-structure}, and its connection one-form is
purely imaginary, so it is metric compatible; uniqueness identifies it
with the Chern connection.  Since $\bar\partial G=0$ and
$\partial\bar G=0$, its curvature is the global $(1,1)$ form
\begin{equation}
 F_{\nabla_z}
 =-2\pi i\bigl(z\,\partial G+
                 \bar z\,\bar\partial\bar G\bigr)
 =-2\pi i\bigl(za+\bar z\,a^\dagger\bigr).
\label{eq:bcpr-unit-chern-curvature}
\end{equation}
With the Bott--Chern convention
\begin{equation}
 c_{1,BC}(L,h)=
 \left[\frac{i}{2\pi}F_{\nabla^h}\right]_{BC}
\label{eq:bcpr-chern-convention}
\end{equation}
of \cite[\S8.a]{Schweitzer2007},
\cref{eq:bcpr-unit-chern-curvature} gives, for every $z\in\mathbb C$,
\begin{equation}
 c_{1,BC}\bigl(\exp(z\mathfrak g)\bigr)
 =z\mathsf A+\bar z\,\mathsf B.
\label{eq:bcpr-chern-exp-identity}
\end{equation}
Consequently, if $\varepsilon=1$ for the orientation fixed by
\cref{eq:bcpr-base-orientation} and $\varepsilon=-1$ for its opposite,
then
\begin{equation}
 c_{1,BC}\bigl(f^*\mathcal O_B(1)\bigr)
 =12\varepsilon\bigl(\mathsf A+\mathsf B\bigr).
\label{eq:bcpr-c1-pullback}
\end{equation}
This is the convention-consistent form of the factor-$12$ identity; no
analytic rescaling of the Picard coordinate changes its integral content.

\subsection{All Bott--Chern generators}

Using \cref{eq:bcpr-global-representatives}, define the following
$d$-closed forms and their classes:
\begin{equation}
\begin{array}{c|c|c}
\text{class}&\text{representative}&\text{bidegree}\\ \hline
\mathsf A&a=\partial G&(1,1)\\
\mathsf B&a^\dagger=\bar\partial\bar G&(1,1)\\
\mathsf C&c=\partial Q&(2,1)\\
\overline{\mathsf C}&c^\dagger=\bar\partial\bar Q&(1,2)\\
\mathsf D&d_0=\partial(G\wedge Q)&(2,2)\\
\overline{\mathsf D}&d_0^\dagger=
 \bar\partial(\bar G\wedge\bar Q)&(2,2)\\
\mathsf E&e=\partial(Q^2)&(3,2)\\
\overline{\mathsf E}&e^\dagger=\bar\partial(\bar Q^2)&(2,3)\\
\mathsf T&\nu,\quad\int_X\nu=1&(3,3).
\end{array}
\label{eq:bcpr-generators}
\end{equation}
Here $\nu$ is any smooth positive top-degree form divided by its integral.
The chosen $Q$ cannot be real: its fibre class is real, but the proof below
shows that $\partial Q$ is nonzero globally, whereas a real
$\bar\partial$-closed $Q$ would also be $\partial$-closed.

\begin{proposition}[nonvanishing and bases]
\label{prop:bcpr-generators-are-bases}
\Verified\ Together with the unit $\mathbf1\in H^{0,0}_{BC}(X)$, the entries
in \cref{eq:bcpr-generators} give bases of all the nonzero Bott--Chern groups
in \cref{eq:bcpr-nonzero-dimensions}:
$\mathsf A,\mathsf B$ form the basis in bidegree $(1,1)$ and
$\mathsf D,\overline{\mathsf D}$ form the basis in bidegree $(2,2)$;
each other listed bidegree is one-dimensional and is generated by its
displayed class.
\end{proposition}

\begin{proof}
Because $X$ is compact and connected, a smooth function annihilated by both
$\partial$ and $\bar\partial$ is constant.  There is no
$\partial\bar\partial$ image in bidegree $(0,0)$, so
$\mathbf1:=[1]_{BC}$ generates
$H^{0,0}_{BC}(X)=\mathbb C\mathbf1$.

The Dolbeault image of $\mathsf A$ is
$d_1\mathfrak g\neq0$ by injectivity in
\cref{eq:bcpr-hodge-row-one}.  Thus $\partial G$ is a nonzero
Bott--Chern class.  The Dolbeault image of $\mathsf B$ is zero because
$a^\dagger=\bar\partial\bar G$, while
$\mathsf B=\overline{\mathsf A}\neq0$.  Hence
$\mathsf A,\mathsf B$ are independent, and
\cref{eq:bcpr-nonzero-dimensions} makes them a basis.

The classes $d_1\mathfrak g$ and $\mathfrak q$ are independent in
$H^{1,1}_{\bar\partial}(X)$: the first restricts to zero on every smooth
torus fibre, while the second restricts to the nonzero class
by \cref{eq:bcpr-q-restriction,eq:bcpr-q-square-integral}.  Exactness of
\cref{eq:bcpr-hodge-row-one} therefore gives
\begin{equation}
 d_1\mathfrak q=[\partial Q]_{\bar\partial}\neq0.
\label{eq:bcpr-c-nonzero}
\end{equation}
Thus $\mathsf C\neq0$, and conjugation gives
$\overline{\mathsf C}\neq0$.

Set
\begin{equation}
 v=[G\wedge Q]_{\bar\partial}
 \in H^{1,2}_{\bar\partial}(X).
\label{eq:bcpr-v-class}
\end{equation}
Its fibre restriction is the class represented by
$\Gamma\wedge Q_F$.  This class is nonzero, rather than merely represented
by a nonzero invariant form: wedging with the holomorphic one-form
$d\zeta_2$ and using \cref{eq:bcpr-q-gamma-product} gives
\begin{equation}
 \int_{F_b}Q_F\wedge\Gamma\wedge d\zeta_2
 =\frac1{4y\Delta}
   \int_{F_b}d\zeta_1\wedge d\zeta_2\wedge
                    d\bar\zeta_1\wedge d\bar\zeta_2
 =-1
\label{eq:bcpr-q-gamma-serre-pairing}
\end{equation}
for the complex orientation: the middle integral equals
$-4y\Delta>0$.  Thus the displayed constant top form has nonzero integral.
Dolbeault--Serre duality on $F_b$ therefore gives
$[\Gamma\wedge Q_F]_{\bar\partial}\ne0$.  Equivalently, if
$Q_F\wedge\Gamma=\bar\partial_F\Psi$ with $\Psi$ of bidegree $(1,1)$,
then, since $\bar\partial_Fd\zeta_2=0$,
\[
 Q_F\wedge\Gamma\wedge d\zeta_2
 =\bar\partial_F(\Psi\wedge d\zeta_2),
\]
and Stokes' theorem would make the integral in
\cref{eq:bcpr-q-gamma-serre-pairing} zero.  Here
$Q_F\wedge\Gamma$ has bidegree $(1,2)$,
$d\zeta_2$ has bidegree $(1,0)$, and
$\Psi\wedge d\zeta_2$ has bidegree $(2,1)$; its $\partial_F$ is zero in
complex dimension two, so the displayed $\bar\partial_F$-boundary is a
$d$-boundary.  Thus $v\neq0$, and it spans
$H^{1,2}_{\bar\partial}(X)$.  Injectivity in
\cref{eq:bcpr-hodge-row-two} gives
\begin{equation}
 d_1v=[\partial(G\wedge Q)]_{\bar\partial}\neq0.
\label{eq:bcpr-d-nonzero}
\end{equation}
Hence $\mathsf D\neq0$, and conjugation gives
$\overline{\mathsf D}\neq0$.  The latter has zero Dolbeault image because
its representative is $\bar\partial$-exact, so it is independent of
$\mathsf D$; \cref{eq:bcpr-nonzero-dimensions} gives the asserted
$(2,2)$ basis.

The class $[Q^2]_{\bar\partial}$ has nonzero fibre restriction by
\cref{eq:bcpr-q-square-integral}.  Since $v$ spans
$H^{1,2}_{\bar\partial}(X)$, the image of $d_1$ is spanned by $d_1v$.
On a smooth fibre,
\[
 i_b^*(d_1v)
 =d_1[\Gamma\wedge Q_F]_{\bar\partial}
 =[\partial_F(\Gamma\wedge Q_F)]_{\bar\partial}=0,
\]
because $\Gamma$ and $Q_F$ have constant fibre coefficients.  Thus every
class in $\operatorname{im}(d_1\colon H^{1,2}_{\bar\partial}\to
H^{2,2}_{\bar\partial})$ restricts to zero.
Therefore $[Q^2]_{\bar\partial}$ is independent of $d_1v$.  Exactness of
\cref{eq:bcpr-hodge-row-two} yields
\begin{equation}
 d_1[Q^2]_{\bar\partial}
 =[\partial(Q^2)]_{\bar\partial}\neq0.
\label{eq:bcpr-e-nonzero}
\end{equation}
Hence $\mathsf E\neq0$, and conjugation gives
$\overline{\mathsf E}\neq0$.  Finally, $\int_X\nu=1$ proves
$\mathsf T\neq0$.  The dimensions in
\cref{eq:bcpr-nonzero-dimensions} complete the proof.
\end{proof}

\subsection{Every Bott--Chern product}

\begin{theorem}[square-zero positive ideal]
\label{thm:bcpr-square-zero-algebra}
\Verified\ Let $N\subset H^{*,*}_{BC}(X)$ be the span of the nine positive-degree
generators in \cref{eq:bcpr-generators}.  Then
\begin{equation}
 H^{*,*}_{BC}(X)=\mathbb C\mathbf1\oplus N,
 \qquad N^2=0.
\label{eq:bcpr-algebra}
\end{equation}
Equivalently, every product of two positive-degree Bott--Chern classes is
zero.
\end{theorem}

\begin{proof}
We first calculate the products landing below top degree.  Since
$\bar\partial G=\bar\partial Q=0$ and, after conjugation,
$\partial\bar G=\partial\bar Q=0$, the graded Leibniz rule gives the exact
identities
\begin{equation}
 \partial\bar\partial(G\wedge\bar G)
 =-a\wedge a^\dagger,
\label{eq:bcpr-product-ab}
\end{equation}
\begin{equation}
 \partial\bar\partial(G\wedge\bar Q)
 =-a\wedge c^\dagger,
 \qquad
 \partial\bar\partial(\bar G\wedge Q)
 =a^\dagger\wedge c.
\label{eq:bcpr-product-acbar-bc}
\end{equation}
Consequently
\begin{equation}
 \mathsf A\mathsf B=0,\qquad
 \mathsf A\overline{\mathsf C}=0,\qquad
 \mathsf B\mathsf C=0.
\label{eq:bcpr-first-zero-products}
\end{equation}

Choose the pullback Fubini--Study metric on
$f^*\mathcal O_B(1)$.  Its Chern form is the pullback of a $(1,1)$ form on
the curve $B$, so its square vanishes pointwise.  Squaring
\cref{eq:bcpr-c1-pullback} and using
\cref{eq:bcpr-first-zero-products} gives
\begin{equation}
 \mathsf A^2+\mathsf B^2=0.
\label{eq:bcpr-sum-of-squares}
\end{equation}
\Cref{eq:bcpr-sum-of-squares} alone does not separate the two
squares.  Apply $\pi_{\bar\partial}^{2,2}$.  The class
$\mathsf B^2$ has zero Dolbeault image because
\begin{equation}
 (a^\dagger)^2
 =\bar\partial(\bar G\wedge a^\dagger).
\label{eq:bcpr-b-square-dolbeault-zero}
\end{equation}
It follows from \cref{eq:bcpr-sum-of-squares} that
$[a^2]_{\bar\partial}=0$.  On the other hand,
\begin{equation}
 [a^2]_{\bar\partial}
 =d_1[G\wedge a]_{\bar\partial},
 \qquad
 \partial(G\wedge a)=a^2.
\label{eq:bcpr-a-square-d1}
\end{equation}
The first arrow in \cref{eq:bcpr-hodge-row-two} is injective, so
\cref{eq:bcpr-a-square-d1} implies
$[G\wedge a]_{\bar\partial}=0$.  There is therefore a form
$\Phi\in\mathcal A^{1,1}(X)$ such that
\begin{equation}
 G\wedge a=\bar\partial\Phi.
\label{eq:bcpr-phi-primitive}
\end{equation}
Applying $\partial$ gives
\begin{equation}
 a^2=\partial\bar\partial\Phi,
\label{eq:bcpr-a-square-bc-zero}
\end{equation}
and hence $\mathsf A^2=0$.  Conjugation gives
$\mathsf B^2=0$.

It remains among the products below top bidegree to prove
$\mathsf A\mathsf C=0$; its conjugate will give
$\mathsf B\overline{\mathsf C}=0$.  First, the natural map
\begin{equation}
 \pi_{\bar\partial}^{3,2}\colon H^{3,2}_{BC}(X)
 \longrightarrow H^{3,2}_{\bar\partial}(X)
\label{eq:bcpr-bc-to-dolbeault-32}
\end{equation}
is an isomorphism.  Indeed, both spaces have dimension one by
\cref{eq:bcpr-nonzero-dimensions,cor:exact-hodge-diamond}, and
\cref{eq:bcpr-e-nonzero} says that $\mathsf E$ has nonzero image.
Pair $\pi_{\bar\partial}^{3,2}(\mathsf A\mathsf C)$ with the generator
$[G]_{\bar\partial}\in H^{0,1}_{\bar\partial}(X)$.  By
\cref{eq:bcpr-phi-primitive}, $\bar\partial c=0$, and Stokes' theorem,
\begin{equation}
 \int_Xa\wedge c\wedge G
 =-\int_XG\wedge a\wedge c
 =-\int_X\bar\partial\Phi\wedge c
 =-\int_X\bar\partial(\Phi\wedge c)=0.
\label{eq:bcpr-serre-pairing-zero}
\end{equation}
Dolbeault--Serre duality makes
\begin{equation}
 H^{3,2}_{\bar\partial}(X)\times
 H^{0,1}_{\bar\partial}(X)\longrightarrow\mathbb C,
 \qquad
 ([\alpha],[\rho])\longmapsto\int_X\alpha\wedge\rho,
\label{eq:bcpr-serre-pairing}
\end{equation}
a perfect pairing.  Since
$H^{0,1}_{\bar\partial}(X)=\mathbb C[G]_{\bar\partial}$,
\cref{eq:bcpr-serre-pairing-zero} implies
$\pi_{\bar\partial}^{3,2}(\mathsf A\mathsf C)=0$; injectivity in
\cref{eq:bcpr-bc-to-dolbeault-32} gives
\begin{equation}
 \mathsf A\mathsf C=0,
 \qquad
 \mathsf B\overline{\mathsf C}=0.
\label{eq:bcpr-second-zero-products}
\end{equation}
This step uses Dolbeault--Serre duality together with the independently
proved injectivity of \cref{eq:bcpr-bc-to-dolbeault-32}.  The
Bott--Chern--Aeppli pairing of
\cite[\S2.c]{Schweitzer2007} by itself would supply a perfect pairing but
would not annihilate this product.

The only products not yet considered which can have bidegree at most
$(3,3)$ are
\begin{equation}
 \mathsf A\mathsf D,\quad
 \mathsf A\overline{\mathsf D},\quad
 \mathsf B\mathsf D,\quad
 \mathsf B\overline{\mathsf D},\quad
 \mathsf C\overline{\mathsf C},
\label{eq:bcpr-top-products}
\end{equation}
all of bidegree $(3,3)$.  Every non-top generator in
\cref{eq:bcpr-generators} is de Rham exact:
\begin{equation}
\begin{gathered}
 a=dG,\qquad a^\dagger=d\bar G,\qquad
 c=dQ,\qquad c^\dagger=d\bar Q,\\
 d_0=d(G\wedge Q),\qquad
 d_0^\dagger=d(\bar G\wedge\bar Q),\qquad
 e=d(Q^2),\qquad e^\dagger=d(\bar Q^2).
\end{gathered}
\label{eq:bcpr-de-rham-primitives}
\end{equation}
Each factor paired with an exact factor is $d$-closed, so every product in
\cref{eq:bcpr-top-products} is $d$-exact and has integral zero.  Integration
induces an isomorphism
\begin{equation}
 H^{3,3}_{BC}(X)\xrightarrow{\ \int_X\ }\mathbb C.
\label{eq:bcpr-top-integration}
\end{equation}
It is well defined by Stokes' theorem, the source has dimension one by
\cref{eq:bcpr-nonzero-dimensions}, and it is nonzero on $\mathsf T$.
Therefore all five products in \cref{eq:bcpr-top-products} vanish in
Bott--Chern cohomology.

All other products of positive-degree generators have $p>3$ or $q>3$ and
vanish already as differential forms.
\Cref{eq:bcpr-first-zero-products,eq:bcpr-a-square-bc-zero,eq:bcpr-second-zero-products,eq:bcpr-top-products,eq:bcpr-top-integration}
exhaust every possible positive-degree product
and prove \cref{eq:bcpr-algebra}.
\end{proof}

\begin{remark}[literal global gauge boundary]
\label{rem:bcpr-global-gauge-boundary}
The Bott--Chern algebra \cref{eq:bcpr-algebra} and the fibre-coordinate
restrictions displayed in
\cref{eq:bcpr-gamma-form,eq:bcpr-u-form,eq:bcpr-q-fibre,eq:bcpr-q-fibre-expanded,eq:bcpr-q-gamma-product,eq:bcpr-q-square}
are complete.  The canonical data do not select a preferred literal global
coefficient atlas for the smooth forms $G$ and $Q$.  Such an atlas would
require explicit smooth local trivializations and horizontal correction
functions through both finite affine fillings and the cusp.  After fixing
$\mathfrak g$ and one coherent-Leray lift $\mathfrak q$, changing those
smooth gauges changes $G$ and $Q$ by Dolbeault coboundaries but changes
none of the cohomology classes, nonvanishing arguments, or products proved
above.  This is a limitation only on a preferred global coefficient
presentation, not on the representatives as cohomology classes or on the
Bott--Chern algebra.
\end{remark}


\begin{auditrecord}[H2: independent replication of the cotangent calculation]
\label{record:h12-direct-images}
\Open\ A further context-isolated AI derivation has independently recovered
the cyclic invariant lattices (11.4), the determinant divisor (11.6), the
residue cancellation (11.17), the scalar test (11.19)--(11.20), and the
splittings
\[
 R^0f_*\Omega_X^1=\mathcal O(-2),\quad
 R^1f_*\Omega_X^1=\mathcal O\oplus\mathcal O(-1)^{\oplus2},\quad
 R^2f_*\Omega_X^1=\mathcal O\oplus\mathcal O(-1).
\]
It also independently gives $h^{1,2}=1$ and $h^{1,1}=2$.  Independent human
specialist replication remains outstanding, so the record stays \Open; this
external-replication boundary is not an unresolved premise of the theorem
above.
\end{auditrecord}
